\documentclass[11pt]{amsart}
\pdfinfoomitdate=1
\pdftrailerid{}
\pdfsuppressptexinfo=15
\usepackage[T1]{fontenc}
\usepackage{lmodern,amsmath,amssymb,amsthm,mathtools}
\usepackage[margin=1.08in]{geometry}
\usepackage[expansion=false]{microtype}
\usepackage{enumitem,booktabs,array,needspace,tikz-cd}
\usepackage[hidelinks,pdfusetitle]{hyperref}
\usepackage[capitalise,noabbrev,nameinlink]{cleveref}
\hypersetup{pdftitle={Projective Hodge lines and ordinary Iitaka subadditivity},pdfauthor={OpenAI},bookmarksdepth=2}
\newcommand{\C}{\mathbf C}
\newcommand{\Q}{\mathbf Q}
\newcommand{\R}{\mathbf R}
\newcommand{\Z}{\mathbf Z}
\newcommand{\N}{\mathbf N}
\newcommand{\cO}{\mathcal O}
\newcommand{\cV}{\mathcal V}
\newcommand{\ddc}{dd^c}
\DeclareMathOperator{\Gr}{Gr}
\DeclareMathOperator{\End}{End}
\DeclareMathOperator{\Hom}{Hom}
\DeclareMathOperator{\Sym}{Sym}
\DeclareMathOperator{\im}{im}
\DeclareMathOperator{\rank}{rank}
\DeclareMathOperator{\rk}{rk}
\DeclareMathOperator{\ord}{ord}
\DeclareMathOperator{\Supp}{Supp}
\DeclareMathOperator{\vol}{vol}
\DeclareMathOperator{\Spec}{Spec}
\DeclareMathOperator{\Gal}{Gal}
\DeclareMathOperator{\Aut}{Aut}
\DeclareMathOperator{\Pic}{Pic}
\DeclareMathOperator{\id}{id}
\newtheorem{theorem}{Theorem}[section]
\newtheorem{lemma}[theorem]{Lemma}
\newtheorem{proposition}[theorem]{Proposition}
\newtheorem{corollary}[theorem]{Corollary}
\theoremstyle{definition}
\newtheorem{definition}[theorem]{Definition}
\theoremstyle{remark}
\newtheorem{remark}[theorem]{Remark}
\newtheorem{example}[theorem]{Example}
\numberwithin{equation}{section}
\setcounter{tocdepth}{1}
\emergencystretch=1em
\allowdisplaybreaks[2]
\title[Projective Hodge lines and ordinary subadditivity]{Projective Hodge lines and ordinary Iitaka subadditivity}
\author{OpenAI}
\date{September 27, 2026}
\begin{document}
\begin{abstract}
We prove the ordinary Iitaka subadditivity conjecture for surjective
projective morphisms with connected fibers between smooth connected
projective varieties over algebraically closed fields of characteristic
zero. If $F$ is the geometric generic fiber of $f:X\to Z$, then
$\kappa(X)\geq\kappa(F)+\kappa(Z)$.
\end{abstract}
\maketitle
\tableofcontents
\section{Introduction}\label{sec:introduction}\label{altord:section}

Let $f:X\to Z$ be a surjective morphism with connected fibres between
smooth projective varieties. The pluricanonical forms on $X$ measure
its birational complexity. Iitaka's subadditivity problem asks whether
the forms on the base and on the geometric generic fibre always
contribute additively to that complexity. More precisely, the Kodaira
dimension $\kappa(X)$ is the maximum dimension of the images of the
complete linear systems $|mK_X|$, for positive integers $m$; it is
$-\infty$ if every such system is empty. If $F$ is the geometric
generic fibre of $f$, the conjecture $C_{n,m}$ asks for
\[
                 \kappa(X)\geq\kappa(F)+\kappa(Z),
                 \qquad n=\dim X,\quad m=\dim Z.
\]
We use $(-\infty)+a=-\infty$ and $\kappa(\mathrm{point})=0$.

The problem belongs to Iitaka's theory of dimensions of divisors
\cite{IitakaOriginal}. Its development has repeatedly connected the
growth of pluricanonical systems with positivity on the base.
Viehweg's weak-positivity method gives addition when the base is of
general type \cite{ViehwegAdditivity}. Kawamata proved addition over
curves and when the general fibre has a good minimal model
\cite{KawamataCurves,KawamataKodaira}; Koll\'ar established the case
of general-type fibres \cite{Kollar87}. The log-general-type-fibre
theorem of Kov\'acs--Patakfalvi and its formulation by Hashizume
provide the addition theorem used below
\cite{KovacsPatakfalvi,Hashizume}.

Birkar proved the conjecture in total dimension at most six
\cite{BirkarIitakaSix}. Using canonical-bundle and nonvanishing
methods, Chang extended this to total dimension at most seven under
the additional assumption $\kappa(X)\geq0$
\cite[Theorem~1.8]{ChangNonvanishing}.
Other advances include Cao--P\u{a}un's theorem for projective klt pairs over abelian varieties,
Hacon--Popa--Schnell's result for bases of maximal Albanese dimension,
and Cao's theorem for projective klt pairs over surfaces
\cite{CaoPaunAbelian,HaconPopaSchnellAbelian,CaoSurfaces}.
These results show how information about the base and positivity of
direct images can control the total pluricanonical system.
Tsuji's preprint states ordinary
subadditivity in Theorem~1.13, while Maehara's preprint states a
broader variation inequality in Section~7, Theorem~8
\cite{TsujiDirectImages,MaeharaIitaka}. Those preprint claims are not
inputs to our proof. The argument here is a projective proof through
a Hodge line and a canonical bundle formula.

\begin{theorem}[Ordinary Iitaka subadditivity]\label{thm:ordinary}\label{altord:ordinary-subadditivity}
Let $k$ be an algebraically closed field of characteristic zero, and
let $f:X\to Z$ be a surjective projective morphism with connected
fibres between smooth connected projective $k$-varieties. If $F$ is
its geometric generic fibre, then
\[
                     \kappa(X)\geq\kappa(F)+\kappa(Z).
\]
\end{theorem}

Thus the ordinary Iitaka conjecture has a positive answer in this
setting. The main step in the proof is a positivity statement for
the entire highest Hodge line of a rationally polarized variation.
We first describe that statement, then explain how the canonical
bundle formula brings it into the subadditivity problem.

Let $Y$ be a smooth connected projective complex variety and $U$ a
dense open subset. A rationally polarized integral pure variation
$\mathbb V$ on $U$ consists of a rational local system with a
monodromy-invariant lattice, a Hodge filtration of fixed weight, and
a rational flat polarization. Suppose its highest nonzero
filtration step $F^p\mathbb V_{\mathcal O}$ is a line, so
$F^{p+1}=0$ and $\operatorname{rank}F^p=1$. On a simply connected
open set, a flat marking identifies this line with a holomorphic
map to the projective space of a fixed reference fibre. We call
this the \emph{highest-line map}. The \emph{full period map}
records every step of the Hodge filtration.

The highest line has a rational parabolic extension across a smooth
normal-crossing compactification. This is the extension whose pullback,
after a finite level cover and resolution making local monodromy
unipotent, is the Deligne--Schmid highest line. The rational weights
retain the finite parts of the original local monodromies.

\begin{theorem}[Projective adjoint positivity]\label{thm:adjoint}\label{altord:p:observation}
Let $Y$ be a smooth connected projective complex variety, let
$U\subset Y$ be a dense Zariski open subset, and let $\mathbb V$
be a rationally polarized integral pure variation of Hodge structure
on $U$ whose entire highest nonzero Hodge filtration step is a line.
Let $M\in\operatorname{Pic}(Y)\otimes\Q$. Assume that, after a smooth
projective birational modification of $Y$ if necessary, there are a
smooth projective alteration $\tau:\widehat Y\to Y$ and a rational
number $c>0$ such that the pulled-back variation has unipotent
monodromy along a simple-normal-crossing boundary and
\[
                     \tau^*M\sim_{\Q}cJ_{\widehat Y},
\]
where $J_{\widehat Y}$ is its extended highest line.
After further smooth projective birational modification, there are a
surjective morphism $p:Y\to S$ with connected fibres onto a smooth
projective variety and a nef rational line $L$ on $S$ such that
\[
 M\sim_{\Q}p^*L,\qquad
 K_S+jL\text{ is big for all sufficiently large integers }j.
\]
If $S$ is a point, the conclusion is $M\sim_{\Q}0$.
\end{theorem}

The two period maps have separate jobs. The full period image provides
the projective base on which an adjoint divisor can be big. The
numerical dimension of $L$ is measured by the highest-line map,
whose rank can be smaller. The argument must therefore remove the
boundary using the rank of the line, without identifying it with
the dimension of the full period image.

\subsection{The boundary obstruction and its resolution}

The algebraicity theorem of Bakker--Brunebarbe--Tsimerman constructs
the full period image as an algebraic variety
\cite[Theorem~1.1]{BakkerBrunebarbeTsimerman}. At a neat level,
Brunebarbe--Cadorel's theorem gives logarithmic general type when
the full period map is generically immersive
\cite[Theorem~1.1]{BrunebarbeCadorel}. A finite quotient returns to
the original level. This introduces a second divisor to be removed:
besides the boundary with nonzero monodromy logarithm, there is
divisorial ramification of the finite quotient. Ordinary adjoint
positivity requires control of both.

We first retain the smallest rational subvariation containing the
highest line. This is the pure transcendental-part construction
of Bakker--Filipazzi--Mauri--Tsimerman
\cite[Definition~2.9]{BFMT}. Rational minimality has a useful
consequence: a connected rational normal subgroup of monodromy
that fixes the line pointwise must be trivial. Together with
Andr\'e's monodromy normality theorem \cite[\S5]{AndreMonodromy}
and the fixed-part results of Griffiths, Deligne, and Schmid
\cite[added note to~4.2.6 and~4.2.8(i),(ii)]{DeligneHodgeII}
\cite[Theorem~7.22 and Corollary~7.23]{Schmid},
it also detects a flat automorphism
that is scalar on all highest lines and preserves one full period.
Such an automorphism preserves every full period.

The local analytic step concerns a family of holomorphic maps
approaching a divisor. If its limiting tangential rank equals
the full rank in the interior, then its nearby image is already
contained in the limiting image. At a nilpotent degeneration,
the limiting line is in the kernel of the monodromy logarithm;
rational minimality rules out equality of ranks. At a divisor
fixed by finite inertia, the limiting lines lie in one eigenspace.
The same rank principle and the fixed-part theorem would make
the inertia element fix the entire full period image.

For this last contradiction we use the normalization of the
\emph{actual} period image. The finite group acts effectively on
its function field. Equality of the full periods therefore forces
the inertia element to be the identity. A further finite cover
of the image could have a nontrivial automorphism acting trivially
on every period, so connected-fibre factorization is taken only
after this argument. Sections~\ref{sec:rational}--\ref{sec:rankloss}
give the rational, analytic, and finite-quotient parts of this rank
loss in that order.

Once the ranks drop, elementary section counting removes the
boundary and divisorial ramification. If the highest-line map has
generic rank $r$, the ample
perturbation in the interior supplies a polynomial of degree $r$
in a large multiple of $L$, whereas the restriction polynomials
on those divisors have degree at most $r-1$. Effective divisors
exceptional over the quotient do not change its section growth.
This yields the adjoint positivity in Theorem~\ref{thm:adjoint}.

\subsection{From the period theorem to ordinary subadditivity}

For the nonvacuous case of Theorem~\ref{thm:ordinary}, put
$d=\kappa(F)\geq0$ and assume $\kappa(Z)\geq0$. The relative Iitaka
fibration factors a birational model of $X$ through a variety $Y$
with $\dim(Y/Z)=d$ and Kodaira-dimension-zero generic fibre over
$Y$. The canonical bundle formula separates the singular-fibre
contribution from a nef moduli divisor:
\[
 K_{\widetilde X}\sim_{\Q}
 g^*(K_Y+B+M)+R,\qquad \widetilde X\xrightarrow{g}Y\xrightarrow{q}Z.
\]
Here $B$ is an effective klt boundary, and the residual divisor $R$
has the direct-image and exceptionality properties needed to compare
complete pluricanonical systems. This is the strategy of
Fujino--Mori's canonical bundle formula, with the birational
discriminant and moduli theory of Ambro
\cite[Theorem~4.5]{FujinoMori}
\cite{AmbroBoundary,AmbroModuli}.

The least-index canonical root cover of the Kodaira-zero fibre
has geometric genus one. Its unique top form is therefore the
entire highest line of a rational variation, even if the character
of the cyclic covering group is not rational. The Hodge realization
of $M$ matches its extension on an alteration, not only its
restriction to the open base. These facts place $M$ within
Theorem~\ref{thm:adjoint}. The auxiliary sub-boundary in this formula
can be negative on the generic fibre; the argument uses nefness
and the Hodge realization, not a semiampleness assertion requiring
an effective generic boundary.

The section comparison gives
\[
 \kappa(X)\geq\kappa(Y,K_Y+B+M),\qquad
 (K_Y+B+M)|_{Y_{\bar\eta}}\text{ big},
\]
where $Y_{\bar\eta}$ is the geometric generic fibre of $q$.
The period projection writes $M=p^*L$ for a morphism $p:Y\to S$.
Write $\mu:\widetilde X\to X$ for the birational modification
of the original source in the relative Iitaka construction.
Figure~\ref{fig:two-bases} separates the two maps from $Y$: their
bases need not map to one another.

\begin{figure}[ht]
\centering
\begin{tikzcd}[column sep=large,row sep=large]
\widetilde X \arrow[r,"g"] \arrow[dr,"f\mu"'] &
Y \arrow[d,"q"] \arrow[r,"p"] & S \\
& Z &
\end{tikzcd}
\caption{The relative Iitaka map $g$ has Kodaira-dimension-zero
generic fibre. The original base map $q$ and the independent period
map $p$ start from the same variety $Y$. The divisor $M=p^*L$
comes from $S$, while the desired addition inequality is measured
over $Z$.}
\label{fig:two-bases}
\end{figure}


A product-map argument first gives log pluricanonical nonvanishing on a general fibre
of $p$. Fujino's twisted weak positivity then turns the big
divisor $K_S+jL-H$, for an ample $H$ and large $j$, into an
effective divisor linearly equivalent to
$K_Y+B+jM-p^*H$ \cite[Theorem~1.1]{FujinoWeakPositivity}.
Interpolating this divisor with an effective klt log canonical
class reduces the desired inequality to the established addition
theorem for log-general-type fibres
\cite[Theorem~2.11]{Hashizume}. Scalar extension of each
pluricanonical space transfers the result from $\C$ to any
algebraically closed field of characteristic zero.

There are two further proofs of the Hodge-line boundary estimate
in the projective setting. An integral tensor replacement makes
monodromy faithful on the moving highest line; projection to a
fixed quotient by a monodromy logarithm then produces transverse
loops on which the line is constant. A second calculation obtains
the restriction-intersection bound directly from Hodge curvature:
smooth Thom forms concentrate on a boundary divisor, and a product
Poincar\'e estimate controls their limits at crossings. These
two calculations share the integral tensor replacement and the
transverse-loop criterion; the second supplies a different proof
of the boundary restriction-intersection step. These
arguments retain the same whole-top-line hypothesis but isolate
different mechanisms behind the boundary estimate. In the main
route, rational minimality and the tangential rank principle treat
the finite quotient on the actual period image directly.

\subsection{Organization}

Section~\ref{sec:rational} constructs the rationally minimal variation and proves the
fixed-period criterion. Section~\ref{sec:limits} compares numerical rank with
boundary limits and proves the tangential rank principle.
Section~\ref{sec:rankloss} applies these results to boundary monodromy and divisorial
ramification. Section~\ref{sec:quotient} constructs the actual period quotient and
completes the proof of adjoint positivity. Section~\ref{sec:canonical} gives the
canonical bundle data, the least-index root cover, and the
section-growth comparison. Section~\ref{sec:ordinary} proves ordinary subadditivity
and gives the scalar-extension argument in characteristic zero.
Section~\ref{p1logalt:section} proves the integral tensor replacement and the alternative
adjoint argument by transverse loops. Section~\ref{p1varhd:section} proves the
restriction-intersection estimate using Thom forms, including the
crossing calculation.

\subsection{Conventions}
An algebraic fibre space is a surjective projective morphism with
connected fibres. All varieties in the Hodge-theoretic argument are
over $\C$. Divisors and line-bundle identities are rational unless
otherwise specified, and section comparisons are taken in degrees
clearing all denominators. We write $A\sim_{\Q}B$ when a positive
integral multiple of $A-B$ is principal. An effective rational class
means a class rationally linearly equivalent to an effective divisor.
The positive and negative parts of a rational divisor have disjoint
support. Numerical dimension is used only for nef divisors and is
defined in Section~\ref{sec:limits}.

\section{Rational minimality and fixed periods}\label{sec:rational}\label{altord:h:section}
Our objective is to recognize when a symmetry that fixes all highest
lines must fix the full periods. We first discard the rational
subvariations that do not contribute to the highest line. Rationality
is essential: the invariant spaces of rational normal monodromy
subgroups will then be Hodge substructures.

Let $U$ be a smooth connected complex algebraic variety and let
$V=(V_{\Z},F^\bullet V_{\mathcal O},Q)$ be a polarized integral pure
variation of Hodge structure of weight $w$. Integral means that the
underlying rational local system has a monodromy-invariant lattice;
the polarization is rational and need not be integral. Write $m$ for the largest index with
$F^mV_{\mathcal O}\ne0$, and assume throughout that
$\rk F^mV_{\mathcal O}=1$. We call this piece the \emph{Hodge line}.
On a simply connected flat chart its inclusion in $V_{\mathcal O}$
defines a holomorphic map
\[
 \ell:U_{\mathrm{loc}}\longrightarrow \mathbb P(V_{\C,u}).
\]
Changes of flat chart compose $\ell$ with a constant projective linear
transformation, so its differential rank is intrinsic.

\subsection{The smallest rational subvariation}

We call a lifted point \emph{Hodge-generic} if it lies outside every
proper Hodge locus of a rational flat tensor in tensor constructions on
the variation and its dual. These are countably many closed analytic
loci on the universal cover. Their complement is dense by the Baire
theorem. Only this analytic genericity is needed in the Hodge arguments.

\begin{lemma}[Rationally minimal variation]\label{altord:h:minimal}
There is a unique smallest rational subvariation $V^{\min}\subset V_{\Q}$
whose Hodge filtration contains $F^mV_{\mathcal O}$. At a Hodge-generic
point, its fibre is the smallest rational Hodge substructure containing
the Hodge line, and that fibre is simple. The subvariation has an induced
polarization and integral lattice. On a smooth compactification with
unipotent boundary monodromy, its extended Hodge line agrees with that
of $V$.
\end{lemma}

\begin{proof}
In a polarizable rational Hodge structure, intersect all rational Hodge
substructures containing the specified line. A finite intersection
already attains the minimum dimension, and gives the unique smallest
one, say $W$. Semisimplicity shows that $W$ is simple: in a decomposition
into simple summands exactly one summand can have a nonzero deepest
Hodge piece, since that piece has dimension one, and that summand already
contains the line.

Here is also the generic-fibre justification, which is useful when
passing between variations and individual Hodge structures. Trivialize
the rational local system on the connected universal cover of $U$.
For each rational endomorphism $e$, the condition that $e$ preserve every
$F^p$ is closed analytic. Choose a lifted point outside all such loci
which are proper. A rational Hodge projector onto $W$ at this point
therefore preserves the filtration everywhere on the cover. Its action
on the Hodge line is the identity everywhere: it is an idempotent on a
line bundle, and equals the identity at the chosen point. Its constant
image is thus a subvariation containing the line. At every other point
chosen in the same way, applying the argument in both directions shows
that its image is again the smallest Hodge substructure containing the
line. Uniqueness and the deck action now imply monodromy invariance,
so this subvariation descends to $U$. Any rational subvariation containing
the line contains it, by inspection at the chosen point.

Restriction of $Q$ to a Hodge substructure is nondegenerate, and
$V^{\min}\cap V_{\Z}$ supplies the lattice. The polarization gives a
complementary subvariation whose deepest Hodge piece is zero.
Canonical extensions preserve this direct sum, proving the last claim.
This is the pure case of the transcendental-part construction of
\cite[Definition~2.9]{BFMT}.
\end{proof}

\subsection{From a fixed highest line to a fixed full period}
A variation is \emph{minimal for its highest line} if the generic
fibre is the smallest rational Hodge substructure containing that
line. Lemma~\ref{altord:h:minimal} supplies this replacement without
changing the line. The next two lemmas explain what this minimality
buys. We use the usual Mumford--Tate group of a rational Hodge
structure: the smallest rational algebraic group through which its
Hodge homomorphism factors.

\begin{lemma}[Normal subgroups and the deepest line]
\label{altord:p:normality}
Let $\mathbb V$ be a rationally polarized integral pure variation on a smooth
connected complex algebraic variety, minimal for its deepest Hodge line.
In a flat trivialization on the universal cover, let $G\subset\mathrm{GL}(V)$
be connected algebraic monodromy and let $P=\operatorname{MT}(V)$ be the generic
Mumford--Tate group.  If $H\subset G$ is a connected normal algebraic
subgroup defined over $\Q$, and $V^H_{\C}$ contains the deepest line at a
Hodge-generic lifted point, then $H$ is trivial.
\end{lemma}

\begin{proof}
The monodromy and fixed-part theorems give that $G$ is semisimple and
normal in the connected reductive group $P$; see
\cite[\S5, Theorem~1 and Corollary~1]{AndreMonodromy}.  Over an
algebraic closure, a connected normal subgroup of a semisimple group is
the product of a collection of its almost simple factors.  Conjugation
by $P$ preserves $G$.  Since $P$ is connected, it acts trivially on the
finite set of these factors, and therefore preserves $H$.

Consequently $V^H$ is a rational $P$-subrepresentation.  It is a rational
Hodge substructure at a Hodge-generic point, and it contains the deepest
Hodge line.  Minimality gives $V^H=V$.  The inclusion
$G\subset\mathrm{GL}(V)$ is faithful, so $H=1$.
\end{proof}

\begin{lemma}[A flat automorphism over a fixed point]
\label{altord:p:inertia}
Let $U$ be a smooth connected complex algebraic variety with a rationally
polarized integral pure variation $\mathbb V$ minimal for its deepest Hodge line.
Suppose that an automorphism $b$ of $U$ fixes a point $e$, and that there
is a flat isomorphism of variations
\[
  \Phi_y:V_y\longrightarrow V_{b(y)}.
\]
Use the fibre $V_e$ as the flat reference space, and put $a=\Phi_e$.
If, on a nonempty open subset of the universal cover, all deepest Hodge
lines lie in a single eigenspace of $a$, then $a$ preserves the complete
Hodge filtration at every lifted point.  In particular, locally at $e$
the full lifted period maps at $y$ and $b(y)$ are equal.
\end{lemma}

\begin{proof}
Write $\rho$ for monodromy based at $e$.  Flatness and $b(e)=e$ give
\[
  a\rho(\gamma)a^{-1}=\rho(b_*\gamma)
  \qquad (\gamma\in\pi_1(U,e)).
\]
Thus $a$ normalizes both algebraic monodromy and its identity component
$G$.  Let $\lambda$ be the eigenvalue in the hypothesis.  The condition
that the deepest line lies in $\ker(a-\lambda)$ is a holomorphic linear
condition on the universal cover, so it holds throughout that connected
cover by analytic continuation.

Fix a Hodge-generic lifted point and let $W\subset V_{\C}$ be the complex span
of the monodromy orbit of its deepest line.  This space is monodromy
invariant, hence $G$-invariant. Equivariance places every line in this
orbit among the deepest lines on the universal cover, so
$a|_W=\lambda\,\mathrm{id}_W$.
Every element $aga^{-1}g^{-1}$, with $g\in G$, therefore acts trivially
on $W$, as do all its $G$-conjugates.  Let $H$ be their normal algebraic
closure in $G$.  It is defined over $\Q$, since $a$ and $G$ are rational.
It is connected: the commutator map from the connected group $G$ has
connected image containing the identity, and the algebraic subgroup
generated by its conjugates is the closure of products of such
connected sets. Although $W$ need not be rational, $V^H$ is rational
and contains the chosen deepest line.
Lemma~\ref{altord:p:normality} gives $H=1$.  Hence $a$ centralizes $G$.

Pass to the connected finite \'etale cover $U_1\to U$ corresponding to
the inverse image of $G$ under monodromy, and choose $e_1$ over $e$.
Since $a$ centralizes all monodromy on $U_1$, it extends to a global flat
rational section of $\End(\mathbb V|_{U_1})$.  The latter is a polarizable
pure variation of weight zero. The flat section $a$ has Hodge type
$(0,0)$ at $e_1$, since $\Phi_e$ is a Hodge isomorphism. The base
$U_1$ is a smooth algebraic variety and hence a Zariski open subset
of a compact complex variety, so the fixed-part theorem applies.
By \cite[Theorem~7.22 and Corollary~7.23]{Schmid}, $a$ has type
$(0,0)$ everywhere. Thus $a$ preserves the full filtration on the
universal cover.

For completeness, take the lift of $b$ to the universal cover fixing a
chosen lift of $e$.  Flatness of $\Phi$ gives
$F^\bullet(b\widetilde y)=aF^\bullet(\widetilde y)$ in the reference
space.  The assertion just proved makes the right-hand side equal to
$F^\bullet(\widetilde y)$.
\end{proof}

\section{Numerical rank and tangential limits}\label{sec:limits}
We now compare a boundary restriction of the extended highest line
with the limiting projective line map. The Hodge-theoretic input
identifies the limiting graded line. The analytic rank principle
then decides what would happen if its tangential rank did not drop.
For a nef line $J$ on a smooth projective $t$-fold we use
\[
 \nu(J)=\max\{a\in\{0,\ldots,t\}:c_1(J)^aH^{t-a}>0\},
\]
where $H$ is ample. A line on a point has numerical dimension zero.

Let $U$ and its variation be as in Section~\ref{sec:rational}.
Take a smooth projective compactification $T$ of $U$ with reduced
simple normal crossings boundary $D$, and suppose that all local
monodromies along $D$ are unipotent. Denote the Deligne extension by
$\overline V$ and the Schmid extension of the Hodge line by
$J=F^m\overline V$. We use the following standard inputs.

\begin{theorem}[Extension and positivity inputs]\label{altord:h:inputs}
Under these hypotheses:
\begin{enumerate}
\item The extended Hodge filtration is by subbundles of $\overline V$.
For a morphism $g:(T',D')\to(T,D)$ of smooth projective varieties
with reduced simple normal crossings boundaries and
$g(T'\setminus D')\subset U$, the extended Hodge line of $g^*V$
is $g^*J$. These extensions commute with direct sums and tensor products.
If the monodromy logarithms of some boundary components vanish, the
variation extends purely across those components, including their
crossings away from the remaining boundary.
\item The line $J$ is nef. It is big precisely when the top
Kodaira--Spencer map
\[
 T_U\longrightarrow
 \Hom(F^mV_{\mathcal O},F^{m-1}V_{\mathcal O}/F^mV_{\mathcal O})
\]
is generically injective, equivalently when $\ell$ is generically
immersive.
\end{enumerate}
The first assertions are the unipotent extension and nilpotent-orbit
theorems of \cite{Schmid,CKS}; functoriality is also recorded in
\cite[Lemmas~2.16 and~5.6]{BFMT}. The second assertion is
\cite[Lemma~6.17]{BakkerBrunebarbeTsimerman}.
\end{theorem}

The equivalence in the second assertion uses Griffiths transversality:
the differential of the line map is the indicated Kodaira--Spencer
map followed by the inclusion of $F^{m-1}/F^m$ in $V_{\mathcal O}/F^m$.

\begin{proposition}[Numerical rank]\label{altord:h:numerical-rank}
If $r=\rk_{\mathrm{gen}}d\ell$, then
\begin{equation}\label{altord:h:eq-numerical-rank}
 \nu(J)=r.
\end{equation}
\end{proposition}

\begin{proof}
Put $t=\dim T$, and fix an ample divisor $H$ on $T$.
For $1\le a\le t$, choose a sufficiently general smooth complete
intersection $C_a\subset T$ of dimension $a$, using a sufficiently
large multiple of $H$. It meets the boundary transversely. At general
points of $C_a\cap U$, the restriction of $d\ell$ to $T_{C_a}$ has
rank $\min(a,r)$. To see the required genericity, at a point where
$d\ell$ has rank $r$, the $a$-planes with this property form a nonempty
open subset of the tangent Grassmannian. A sufficiently ample linear
system prescribes the required first-order tangent conditions; Bertini
then gives the assertion for general complete intersections.
By Theorem~\ref{altord:h:inputs}, $J|_{C_a}$ is big exactly when $a\le r$.
Since it is nef, this is equivalent to
$(J|_{C_a})^a>0$, and hence to $c_1(J)^aH^{t-a}>0$.
This proves the formula, also when $r=0$; dimension zero is immediate.
\end{proof}

\subsection{The limiting line along a divisor}

We recall the boundary input in the form needed here. Let $E$ be a
component of $D$, let
$E^\circ=E\setminus\Supp(D-E)$, and let $N$ be its monodromy logarithm.
The monodromy weight filtration $W(N)$ is centred at $w$.
The limiting filtration induces polarized pure variations on its graded
pieces over $E^\circ$, with unipotent monodromy at the remaining boundary.
If $k$ is the unique weight carrying the deepest line, then $J|_E$ is
the Schmid extension of the top line of $\Gr^{W(N)}_k$.
These statements include compatibility at deeper strata;
see \cite[\S2.5.2--\S2.5.3 and Lemma~2.16(3)]{BFMT}.
For a single $N$, the relevant primitive summand is polarized by the
form induced by $Q(-,N^{k-w}-)$ when $k\ge w$.

The local expression below is Schmid's nilpotent-orbit normal form
\cite[Theorem~4.12]{Schmid}; its compatibility at crossings uses
\cite{CKS}. The rank comparison is the consequence needed here.

\begin{proposition}[Boundary line]\label{altord:h:boundary-line}
On a transverse coordinate chart $(u,v)$ about a general point of $E$,
with $E=(u=0)$, write a generator of the lifted Hodge line as
\begin{equation}\label{altord:h:nilpotent-expression}
 e^{zN}a(u,v),\qquad z=\frac{\log u}{2\pi i},
\end{equation}
where $a$ is holomorphic across $u=0$ and $a(0,v)\ne0$.
After shrinking the chart there is an integer $\ell\ge0$ such that
\begin{equation}\label{altord:h:primitive-line}
 N^{\ell+1}a(0,v)=0,\qquad N^\ell a(0,v)\ne0,
\end{equation}
and the line belongs to weight $w+\ell$ in the limiting mixed Hodge
structure. Moreover,
\begin{equation}\label{altord:h:boundary-rank}
 \nu(J|_E)\le
 \rk_{\mathrm{gen}}d\bigl(v\longmapsto[N^\ell a(0,v)]\bigr).
\end{equation}
The map on the right is formed in the chosen flat coordinates.
If $N=0$, take $\ell=0$; the variation then extends purely near the
general point of $E$.
\end{proposition}

\begin{proof}
Put $a_0=a(0,v)$. In the Deligne bigrading of the limiting mixed
Hodge structure, $F^m$ is a single one-dimensional summand $I^{m,q}$:
all first indices are at most $m$, and $\dim F^m=1$.
Set $\ell=m+q-w$. The primitive decomposition on weight $w+\ell$
expresses each summand as a sum of terms
\[
 N^j P_{w+\ell+2j},\qquad j\ge\max(0,-\ell),
\]
where $P_{w+b}$ denotes the primitive part for $b\ge0$.
A contribution to Hodge type $(m,q)$ with $j>0$ would come from
type $(m+j,q+j)$, impossible because its first index exceeds $m$.
Thus $j=0$, $\ell\ge0$, and the top line is primitive.
Hard Lefschetz for the monodromy weight filtration gives nonvanishing
of its image under $N^\ell$ and vanishing under $N^{\ell+1}$ on the
associated graded. Since $N$ has bidegree $(-1,-1)$ in the Deligne
bigrading, these statements hold for $a_0$ itself: a vector in a single
Deligne summand maps injectively to its weight-graded piece.
The integer $\ell$ is constant on a sufficiently small open stratum,
so the identities hold throughout the chosen $v$-chart.

Let $J_E^{\mathrm{gr}}$ be the extended top line of the pure variation
on $\Gr^{W(N)}_{w+\ell}$ along $E^\circ$.
The boundary extension input identifies it with $J|_E$.
Projecting $N^\ell a_0$ to $\Gr^{W(N)}_{w-\ell}$ yields the image
of this graded top line under the flat isomorphism
\[
 N^\ell:\Gr^{W(N)}_{w+\ell}\longrightarrow
          \Gr^{W(N)}_{w-\ell}.
\]
Consequently the projective map of the graded line has rank no greater
than $v\mapsto[N^\ell a_0(v)]$: projection is a fixed linear map on
the flat chart, defined near the line in question.
Apply Proposition~\ref{altord:h:numerical-rank} to the graded variation on
$E^\circ$ and its smooth projective compactification $E$.
This proves \eqref{altord:h:boundary-rank}.
\end{proof}

\subsection{A local rank principle}

The boundary-line proposition has reduced the numerical question to a
projective differential. We next prove the local assertion that will
force a symmetry to fix the interior whenever the limiting rank is
as large as the interior rank. This analytic fact requires control of tangential derivatives
only. A sector below has the form
$S_\epsilon=\{0<|u|<\epsilon,\ \alpha<\arg u<\beta\}$ with a fixed
branch of the argument and finite opening $\beta-\alpha<2\pi$.

\begin{lemma}[Rank and a tangential limit]\label{altord:h:rank-principle}
Let $B\subset\C^n$ be a polydisc and let
$f:S_\epsilon\times B\to M$ be a holomorphic map to a complex
manifold. Suppose $f(u,\cdot)$ converges uniformly on compact subsets
of $B$, as $u\to0$ throughout the sector, to a holomorphic map $f_0$.
Assume that the full generic differential rank of $f$ is at most $r$
and that $\rk df_0(v_0)=r$ at some $v_0\in B$.
Then, after shrinking the sector and taking neighborhoods
$v_0\in B'\Subset B_0\subset B$,
\[
 f(S_{\epsilon'}\times B')\subset f_0(B_0),
\]
where the indicated local image of $f_0$ is an embedded complex
submanifold of dimension $r$. In particular, every fixed closed analytic
subset containing this local image also contains the nearby image of $f$.
\end{lemma}

\begin{proof}
Use a common coordinate chart in $M$, shrinking $B$ and the sector.
Cauchy estimates give convergence of all derivatives in the $B$
variables on smaller polydiscs. The $(r+1)$-minors of $df$ vanish
identically, so those of $df_0$ vanish as well. Thus $f_0$ has constant
rank $r$ near $v_0$. The holomorphic constant rank theorem supplies
source coordinates $v=(x,y)\in\C^r\times\C^{n-r}$ and target
coordinates in which $f_0(x,y)=(x,0)$.

Write $f=(F,G)$ in these coordinates, with
$F(u,x,y)=x+E(u,x,y)$. On fixed smaller polydiscs, both $E$ and
$\partial_xE$ tend uniformly to zero. Choose nested $x$-balls, a
smaller ball $W$ for $w$, and a ball $Y$ for $y$. For sufficiently
small $\epsilon'$, the maps
\[
 x\longmapsto w-E(u,x,y)
\]
send a fixed closed $x$-ball into itself and have Lipschitz constant
less than $1/2$, uniformly for
$(u,w,y)\in S_{\epsilon'}\times W\times Y$.
The contraction theorem gives a unique solution
$x=\chi(u,w,y)$ of $F(u,x,y)=w$ on this common product.
The holomorphic inverse function theorem and uniqueness make $\chi$
holomorphic in all interior parameters; moreover $\chi\to w$
uniformly on smaller products.

In the new source coordinates the map is
\[
 f(u,\chi(u,w,y),y)=(w,H(u,w,y)).
\]
The $w$-derivatives already have rank $r$.
The rank bound therefore forces $\partial_uH=0$ and
$\partial_yH=0$: either derivative, if nonzero, would give an
additional image vector with zero first $r$ coordinates.
Since the sector and $Y$ are connected, $H$ depends only on $w$.
Its limit as $u\to0$ is zero, hence $H=0$.
Shrinking the original source neighborhood so that its $F$-values
belong to $W$ proves the assertion. If $r=0$, the rank bound directly
makes $f$ constant on the connected source, and the same conclusion
follows.
\end{proof}

\begin{lemma}[Nilpotent tangential convergence]\label{altord:h:nilpotent-limit}
Let $N$ be a nilpotent endomorphism of a finite-dimensional complex
vector space, and let $a(u,v)$ be holomorphic on $\Delta\times B$.
Suppose, for a fixed $\ell\ge0$, that
$N^{\ell+1}a(0,v)=0$ and $N^\ell a(0,v)\ne0$ throughout $B$.
On a sector with the above choice of $z=\log u/(2\pi i)$,
\begin{equation}\label{altord:h:eq-nilpotent-limit}
 [e^{zN}a(u,v)]\longrightarrow[N^\ell a(0,v)]
\end{equation}
normally on sufficiently small $v$-charts, with all tangential
derivatives. Thus Lemma~\ref{altord:h:rank-principle} applies whenever the
rank of this limit equals an upper bound for the full rank of the
line map. Its nearby image then lies in $\mathbb P(\ker N)$.
\end{lemma}

\begin{proof}
Choose $k$ with $N^{k+1}=0$, and write
$a(u,v)=a_0(v)+u q(u,v)$. Set $b(v)=N^\ell a_0(v)/\ell!$.
Direct expansion gives
\begin{align*}
 z^{-\ell}e^{zN}a(u,v)
 &=b(v)+\sum_{j<\ell}\frac{z^{j-\ell}}{j!}N^ja_0(v)
       +u\sum_{j=0}^k\frac{z^{j-\ell}}{j!}N^jq(u,v).
\end{align*}
For every fixed tangential derivative, uniformly on compact subsets
of $B$, the remainder is
\[
 O(|z|^{-1})+O(|u|\,|z|^{k-\ell});
\]
the first term is absent when $\ell=0$.
On the sector $|z|$ is comparable to $|\log|u||$, so both terms tend
to zero. Choose a linear functional nonzero on $b(v_0)$ and shrink
the $v$-chart. Dividing by that functional gives convergence with
all tangential derivatives in a fixed affine projective chart,
proving \eqref{altord:h:eq-nilpotent-limit}. Finally $Nb=0$, so the limiting
image belongs to the fixed projective linear subspace
$\mathbb P(\ker N)$, and the last assertion follows from
Lemma~\ref{altord:h:rank-principle}.
\end{proof}

\begin{remark}\label{altord:h:scaling}
All numerical-rank statements are unchanged on replacing a Hodge line
by a positive rational multiple of its class in $\Pic(T)\otimes\Q$.
The local statements also apply with $N=0$, when the line map extends
holomorphically across the divisor.
\end{remark}

\section{Rank loss on the boundary and the ramification divisor}
\label{sec:rankloss}

The preceding lemmas apply to a smooth model of a full period image.
At a neat level the variation is unipotent at the boundary. Returning
to the original level also introduces ramification divisors. We now
prove the strict numerical rank drop for both kinds of divisor.

We specify the finite-group data because the group acting on the
variation can be larger than the effective group acting on its base.
Let $T$ be a smooth connected projective complex variety of positive
dimension, let $D_T$ be a reduced simple-normal-crossing divisor, and
put $U=T\setminus D_T$. Let $\mathbb V$ be a rationally polarized
integral pure variation on $U$, minimal for its entire highest
Hodge line, with unipotent boundary monodromy. Write $J$ for the
extended highest line. Let $D_N$ be the union of the components
of $D_T$ with nonzero monodromy logarithm, and put
$U_N=T\setminus D_N$. The variation extends purely over $U_N$ by
Theorem~\ref{altord:h:inputs}.

Suppose a finite group $H$ acts on $T$, preserves $D_T$, and acts on
the variation by flat Hodge isomorphisms preserving its integral
lattice and rational polarization and covering that action. Set
\[
                 G_f=H/\ker\bigl(H\longrightarrow\Aut(T)\bigr).
\]
Thus $G_f$ acts effectively on $T$. We retain $H$ for its action on
the variation, since its kernel on $T$ can act nontrivially on fibres.
Uniqueness of the pure extension preserves this action over $U_N$.

Fix a connected period domain $\mathcal D$ for the rational
polarized reference space and a neat arithmetic group $\Gamma'$ of
lattice isometries containing the monodromy. Let
\[
                    \varphi:U_N\longrightarrow\Gamma'\backslash\mathcal D
\]
be the full period map. Assume that the normalization $B$ of the
closure of its actual image is an algebraic variety, that
the natural rational map $T\dashrightarrow B$ is birational, and
that $\varphi$ is generically immersive. In particular,
\[
                         \C(T)=\C(B).
\]
These assumptions describe the actual period-image model. The
algebraicity of $B$ is the period-image theorem of
\cite[Theorem~1.1]{BakkerBrunebarbeTsimerman}; here its construction
is an explicit input to the rank-loss statement.

\begin{theorem}[Boundary and ramification rank loss]\label{thm:rankloss}
For the data just specified, put $r=\nu(J)$.
Then $r>0$. If $E$ is either a component of $D_N$, or a component
of $D_T$ fixed pointwise by a nonidentity element of
$G_f$, then
\[
                           \nu(J|_E)<r.
\]
In particular, after placing the divisorial fixed loci in $D_T$
on an equivariant smooth model, the conclusion applies
to every divisorial ramification component of $T\to T/G_f$.
\end{theorem}

\begin{proof}
Proposition~\ref{altord:h:numerical-rank} identifies $r$ with the
generic rank of the highest-line map. We first show that this
rank is positive, then consider the nilpotent and finite cases.

If $r=0$, the highest line is constant on the connected universal
cover of $U_N$. Monodromy preserves it. Its connected algebraic
monodromy group $G$ is semisimple and hence acts trivially on that
line. Lemma~\ref{altord:p:normality}, applied with the normal
subgroup equal to $G$, gives $G=1$. After a finite \'etale cover the
monodromy is trivial, and the fixed-part theorem makes the entire
variation constant. Its full period map then has rank zero,
contrary to generic immersivity on the positive-dimensional variety
$T$. Thus $r>0$.

\smallskip\noindent
\emph{A component with nonzero monodromy logarithm.}
Let $E\subset D_N$ and write $N\ne0$ for its rational monodromy
logarithm. In coordinates $(u,v)$ near a general point of $E$,
with $E=(u=0)$, the lifted highest-line map has the form
\[
                [e^{zN}a(u,v)],\qquad z=\frac{\log u}{2\pi i}.
\]
Proposition~\ref{altord:h:boundary-line} provides an integer
$\ell\geq0$ with
\[
 N^{\ell+1}a(0,v)=0,\qquad N^\ell a(0,v)\ne0,
 \qquad
 \nu(J|_E)\leq\rk_{\rm gen}d\bigl(v\mapsto[N^\ell a(0,v)]\bigr).
\]
Lemma~\ref{altord:h:nilpotent-limit} gives its normal
tangential convergence on a sector from the interior line maps.
All interior $(r+1)$-minors vanish. Cauchy convergence of tangential
derivatives therefore bounds the rank of the limiting map by $r$.

Suppose that $\nu(J|_E)\geq r$. The limiting map has rank exactly
$r$ at a suitable general point, and its image lies in
$\mathbb P(\ker N)$. Lemma~\ref{altord:h:rank-principle} forces all
nearby interior lines to lie in this same linear subspace. The
condition that $N$ annihilates the highest line is holomorphic, so
it continues over the connected universal cover of $U_N$.

The one-parameter group $\exp(tN)$ belongs to $G$: it is the
connected Zariski closure of the cyclic unipotent local monodromy.
Every monodromy conjugate of this group fixes the highest line
pointwise. Let $G_N$ be the algebraic subgroup they generate.
It is connected, defined over $\Q$, and normalized by monodromy;
Zariski density makes it normal in $G$. Its invariant space contains
the highest line at a Hodge-generic point. By
Lemma~\ref{altord:p:normality}, $G_N=1$, contradicting $N\ne0$.
This proves the required inequality for components of $D_N$.

We have controlled all divisors at which the pure variation fails
to extend. The remaining obstruction is a divisor of finite
ramification where the variation does extend. Its inertia element
can act on the rational fibre even when it fixes every point of
the divisor, so we must use the full-period conclusion of
Lemma~\ref{altord:p:inertia}.

\smallskip\noindent
\emph{A divisor fixed by finite inertia.}
Let $b\in G_f$ be nonidentity and fix $E$ pointwise. If $E$ is
contained in $D_N$, the preceding case applies. Otherwise a general
point $e\in E$ lies in $U_N$. Choose a lift $h\in H$ of $b$.
The retained $H$-action supplies a flat Hodge isomorphism
\[
                  \Phi_y:V_y\longrightarrow V_{b(y)}
\]
on $U_N$. Use $V_e$ as a flat reference space on a small simply
connected $b$-invariant neighborhood, and set $a=\Phi_e$.
If $x$ is the full lifted period map, flatness gives
\[
                             x(by)=a\,x(y).
\]
The operator $a$ is a rational Hodge isometry at $e$. Because $H$
acts on the variation and $b(e)=e$, the group-action law gives
$a^{\operatorname{ord}(h)}=1$. Since $b$ fixes $E$, every highest line
on the nearby part of $E$ is an eigenline of $a$. Its eigenvalue
is constant on a connected small patch, so these lines lie in one
eigenspace $\ker(a-\lambda)$.

Proposition~\ref{altord:h:boundary-line} with $N=0$ bounds
$\nu(J|_E)$ by the rank of this extending line map on $E$.

Suppose again that $\nu(J|_E)\geq r$. The rank on $E$ is at most
$r$, since all $(r+1)$-minors of the ordinary extending line map
vanish by continuation from the interior. It is therefore $r$.
Apply Lemma~\ref{altord:h:rank-principle} in a transverse coordinate
with $N=0$. A nonempty open set of lifted highest lines lies in
$\ker(a-\lambda)$. Lemma~\ref{altord:p:inertia}, applied on the
smooth algebraic variety $U_N$ with fixed point $e$, now shows
that $a$ preserves every full period. Consequently $x(by)=x(y)$.

Projecting to $\Gamma'\backslash\mathcal D$ gives
$\varphi(by)=\varphi(y)$ on a nonempty open set. Since $B$ is the
normalization of the actual image, its normalization map is
birational and hence injective on a dense open subset of that
image. It follows that $b$ acts trivially on $\C(B)$. The chosen
model has $\C(T)=\C(B)$, so $b$ acts trivially on $T$, contrary
to its nonidentity in $G_f$. This proves the strict inequality.

Finally, in a finite quotient in characteristic zero, ramification
at a divisorial valuation is detected by a nontrivial inertia
group, which acts identically on that divisor's function field.
Its elements therefore fix the divisor pointwise. The last
assertion follows from the case just proved.
\end{proof}

\begin{remark}
The conclusion concerns the numerical dimension of the highest line,
while generic immersivity is a hypothesis on the full period map.
Neither the proof nor the conclusion identifies these ranks. The
function-field equality with the actual period image is used only
at the last contradiction in the finite case; it cannot be replaced
there by an arbitrary finite extension of that function field.
\end{remark}

\section{The period quotient and adjoint bigness}\label{sec:quotient}\label{altord:sec:period}
We prove Theorem~\ref{thm:adjoint}. The full period image supplies
an algebraic base, and Section~\ref{sec:rankloss} controls the divisors
introduced by compactification and finite descent. We first show
that those numerical rank drops suffice to remove the divisors
from an adjoint class. This is a section-growth calculation; no
abundance statement for the highest line is needed.

\subsection{Section estimates}\label{altord:sec:estimates}
\begin{lemma}\label{altord:e:volume-lemma}
Let \(T\) be a smooth projective \(t\)-fold, \(A\) a rational divisor,
and \(J\) a nef rational divisor
with \(r=\nu(J)>0\), and \(D=\sum_i d_iE_i\) an effective integral divisor
with smooth prime components. Suppose
\[
 \nu(J|_{E_i})<r\quad\text{for every }i.
\]
If \(A+D\) is big, then \(A+jJ\) is big for every sufficiently large integer
\(j\).
\end{lemma}

\begin{proof}
Choose an ample rational divisor \(H_1\) such that \(A+D-H_1\) is
effective. In this proof \(m\) tends to infinity through sufficiently
divisible positive integers. For each fixed \(j\geq0\), asymptotic Riemann--Roch for the ample
divisor \(H_1+jJ\) gives
\[
 \liminf_{m\to\infty}
 \frac{t!}{m^t}h^0\bigl(T,m(A+D+jJ)\bigr)
 \ \geq\ (H_1+jJ)^t\ \geq c j^r
\]
for a constant \(c>0\) and all \(j\geq1\).
The last inequality follows from \(J^rH_1^{t-r}>0\).

Remove \(mD\) one prime divisor at a time. Each quotient in a divisor
exact sequence is a line bundle on some \(E_i\), of the form
\[
 \mathcal O_{E_i}
 \left(m(A+D+jJ)-\sum_h k_hE_h\right),
 \qquad 0\leq k_h\leq md_h.
\]
There is a single ample rational divisor \(H_2\) on \(T\) for which every
such quotient has at most
\(h^0(E_i,m(H_2+jJ)|_{E_i})\) sections, in divisible degrees and for every
\(j\geq0\). To see this, choose an ample Cartier divisor \(B_0\) with a
globally generated multiple clearing the denominator of
\(B_0-(A+D)\), and ample Cartier divisors \(B_h\) such that both \(B_h\)
and \(B_h+E_h\) are globally generated. Enlarge \(B_0\) if necessary and
take \(H_2=B_0+\sum_h d_hB_h\). The difference, restricted to \(E_i\), is
\[
 m(B_0-A-D)+
 \sum_h\bigl((md_h-k_h)B_h+k_h(B_h+E_h)\bigr).
\]
It has a section nonzero on \(E_i\), by global generation and
multiplication, giving the claimed bound.

There are \(m\sum_i d_i\) quotient terms. For fixed \(j\),
asymptotic Riemann--Roch on each \(E_i\) bounds their total contribution on
the scale \(m^t/t!\) by
\[
 t\sum_i d_i\,(H_2+jJ)^{t-1}\cdot E_i.
\]
This is \(O(j^{r-1})\), since every intersection containing at least
\(r\) factors of \(J|_{E_i}\) vanishes. Consequently
\[
 \liminf_{m\to\infty}
 \frac{t!}{m^t}h^0\bigl(T,m(A+jJ)\bigr)
 \ \geq\ c j^r-O(j^{r-1})>0
\]
for all sufficiently large \(j\), proving bigness.
\end{proof}

\begin{lemma}\label{altord:e:exceptional-pullback}
Let \(\pi:T\to S\) be a surjective generically finite morphism of smooth
projective varieties. Let \(E\geq0\) be a rational divisor on \(T\) whose
components have image of codimension at least two in \(S\). For a rational
divisor \(D\) on \(S\), the divisor \(\pi^*D+E\) is big if and only if
\(D\) is big.
\end{lemma}

\begin{proof}
Factor \(\pi\) as \(T\xrightarrow{\rho}\bar S\xrightarrow{\tau}S\),
where \(\bar S\) is the normalization of \(S\) in \(\C(T)\).
Then \(\rho\) is birational and \(\tau\) is finite. Every component of
\(E\) is \(\rho\)-exceptional, since a finite map preserves dimensions.
Normality gives
\(\rho_*\mathcal O_T(mE)=\mathcal O_{\bar S}\) in divisible degrees:
a rational function with poles only on exceptional divisors has no
codimension-one pole downstairs. Projection formula identifies the
sections of \(m(\pi^*D+E)\) with those of \(m\tau^*D\).
Finally, bigness is invariant under surjective finite pullback.
\end{proof}


\subsection{Construction and descent}
\begin{proof}[Proof of Theorem~\ref{altord:p:observation}]
We may first make the birational modification allowed in the hypothesis.
Replace $\mathbb V$ by the minimal rational subvariation from
Lemma~\ref{altord:h:minimal}, and equip it with the intersection lattice and
the restricted polarization.  Its deepest Hodge line and the line
$J_{\widehat Y}$ are unchanged.  We continue to denote this variation
by $\mathbb V$.

\medskip\noindent
\textbf{The actual period image and its finite quotient.}
Choose a rational reference space with lattice and polarization, and
let $\mathcal D$ be the connected period-domain component containing a
chosen lift of the periods. Let $\Gamma$ be the arithmetic group of
integral polarizing isometries preserving this component; it contains
the monodromy of $\mathbb V$. Let $\Gamma'\triangleleft\Gamma$ be a
normal neat subgroup of finite index.  The kernel of
\[
  \pi_1(U)\longrightarrow\Gamma/\Gamma'
\]
defines a connected finite \'etale Galois cover $U'\to U$.  Denote its
deck group by $H$; thus $H$ is the image of this homomorphism, not
necessarily all of $\Gamma/\Gamma'$.

The period map of $\mathbb V|_{U'}$ factors algebraically through the
normalization $B$ of the closure of its image in
$\Gamma'\backslash\mathcal D$.  The space $B$ is a normal
quasiprojective variety, and on a smooth dense open subset its period
realization has generically immersive differential; this is the
period-image theorem \cite[Theorem~1.1]{BakkerBrunebarbeTsimerman}.  Here and below the function
field of $B$ is the function field of this normalized \emph{actual
image}.  We do not replace it by its relative algebraic closure in
$\C(U')$.

Neatness makes the arithmetic action free.  The universal local system
on the quotient is
\[
  (\mathcal D\times V_{\Z})/\Gamma',
\]
with its tautological filtration. On the smooth open of $B$ this
filtration is horizontal: the dominant map $U'\to B$ is generically
submersive and its pullback is the original variation, so Griffiths
transversality holds on a dense open and hence everywhere. Thus the
restricted local system and filtration form a polarizable integral
variation. Moreover, the action of
$\Gamma/\Gamma'$ on this local system is an honest action:
\[
  [x,v]\longmapsto[\gamma x,\gamma v].
\]
Normality of $\Gamma'$ makes the formula independent of the
representative $\gamma$.  Restricting gives an action of $H$ on the
variation over $B$, covering its action on $B$. The action on $B$ is
algebraic: apply uniqueness in \cite[Theorem~1.1]{BakkerBrunebarbeTsimerman} to the period map
and its expression obtained by a deck transformation on $U'$ and the
corresponding arithmetic translation on the target, and then lift to
the normalization. Let
\[
  G_f=H/\ker(H\longrightarrow\operatorname{Aut}(B))
\]
be the effective group on the image.  The full group $H$ will still be
retained for its action on the variation.

Choose a $G_f$-equivariant smooth projective compactification $S'$ of a
smooth dense invariant open of $B$, with snc boundary, by equivariant
compactification and resolution.  All boundary monodromies of its
variation are unipotent.  Indeed they are quasi-unipotent by the
monodromy theorem, and they belong to the neat group $\Gamma'$; their
root-of-unity eigenvalues must therefore all be one.  Let $J'$ denote
the extended deepest Hodge line and put $L'=cJ'$.  This is nef by
Theorem~\ref{altord:h:inputs}.  The $H$-action extends to $J'$ by
functoriality of the unipotent Deligne and Hodge extensions.

Let $q:S'\to Q_0=S'/G_f$ be the finite quotient.  Although an element
of $\ker(H\to G_f)$ may act nontrivially on $J'$, it acts on its fibres
through finite-order characters.  Taking a tensor power divisible by
$|H|$ kills the fibre action of every stabilizer in $H$, including this
kernel.  That tensor power therefore descends to an actual line bundle
on $Q_0$ by finite-group descent.  Consequently there is a rational
line bundle $\overline L$ on $Q_0$ with
\[
  q^*\overline L\sim_{\Q}L'.
\]
It is nef, since its pullback by the finite surjective map $q$ is nef.
Choose a smooth projective resolution $s:S_0\to Q_0$ and put
$L_0=s^*\overline L$.

The equivariant period map on $U'$ descends to a dominant rational map
$Y\dashrightarrow Q_0$, and hence to $S_0$.  Resolve its graph by a
smooth projective birational modification of $Y$, obtaining
\[
  p_0:Y\longrightarrow S_0.
\]
We claim that
\begin{equation}
  M\sim_{\Q}p_0^*L_0.
  \label{altord:p:line-descent}
\end{equation}
To verify this claim, choose a common smooth projective alteration $W$
dominating the alteration in the hypothesis and the level cover, and
resolve the graph of its rational map to $S'$.  Write
$\mu:W\to Y$ and $t:W\to S'$ for the resulting morphisms.  Arrange
snc boundaries by further resolution, containing the inverse images of
the boundaries in question.  The pulled-back variations agree on a
dense open subset.  Unipotent Deligne extensions and their extended
Hodge filtrations commute with these pullbacks between smooth snc
pairs; see \cite[Lemma~5.6]{BFMT}, or the unipotent extension theory
\cite{Schmid,CKS}.  The extended top line $J_W$ is consequently the
pullback of the top line on either model.  The hypothesis and the
commutative maps to $Q_0$ give
\[
  \mu^*M\sim_{\Q}cJ_W
  \sim_{\Q}t^*L'
  \sim_{\Q}\mu^*p_0^*L_0.
\]
Pullback by a proper generically finite map is injective on rational
line-bundle classes: applying divisor pushforward, or the norm of a
principal rational function, multiplies the original class by the
degree.  This proves \eqref{altord:p:line-descent}.  In particular no boundary
twist has been inferred from an equality merely on the open locus.

If $\dim S_0=0$, equation~\eqref{altord:p:line-descent} gives
$M\sim_{\Q}0$, and the theorem follows.  Assume henceforth that
$\dim S_0>0$.

\medskip\noindent
\textbf{A common equivariant model and log general type.}
Normalize $S_0$ in $\C(S')$, and take a $G_f$-equivariant smooth
projective resolution of the resulting model and its birational map
to $S'$.  Denote the resulting maps by
\[
  \eta:T\longrightarrow S',\qquad
  \pi:T\longrightarrow S_0.
\]
Thus $\eta$ is birational, $\pi$ is generically finite with Galois
function-field group $G_f$, and $G_f$ acts on $T$ over $S_0$.
Choose the resolution so that there is a $G_f$-invariant snc boundary
$D_T$ containing the complement of the open period locus and the
support of the Jacobian divisor of $\pi$.  The pulled-back variation
has unipotent boundary monodromy.  If $J_T$ is its extended top line,
then
\begin{equation}
  L_T:=cJ_T\sim_{\Q}\eta^*L'
     \sim_{\Q}\pi^*L_0.
  \label{altord:p:T-line}
\end{equation}

Let $D_N$ be the reduced sum of those components of $D_T$ having
nonzero monodromy logarithm.  The variation extends as a pure polarized
variation over
\[
  U_N=T\setminus D_N.
\]
Indeed, near a point where all the remaining logarithms are zero, the
nilpotent-orbit extension is a pure polarized Hodge filtration; the
several-variable extension gives the same statement at crossings
\cite{Schmid,CKS}.  These extensions are unique and preserve the
$H$-equivariance of the variation.  The divisor $D_N$ is snc, since it
is a union of components of $D_T$.

The full period map on $U_N$ has generically immersive differential,
because $T$ is birational to $B$.  The log-general-type theorem for a
polarized variation with generically immersive period map therefore
gives
\begin{equation}
  K_T+D_N\quad\text{big}.
  \label{altord:p:log-big}
\end{equation}
Here the applicable theorem is \cite[Theorem~1.1]{BrunebarbeCadorel}.  Its boundary is precisely
$D_N$, because the pure variation has just been extended over all the
other boundary components.

We now have precisely the geometric data of
Theorem~\ref{thm:rankloss}: $T$ is birational to the actual normalized
period image, the full period map is generically immersive, and the
retained finite group $H$ acts on its integral polarized variation.
The effective group on $T$ is $G_f$. Minimality is preserved here:
the variation on $T$ and the original variation have the same
Hodge structures at Hodge-generic points, since the original source
dominates the period image. Consequently
\begin{equation}
 r:=\nu(L_T)>0,
 \qquad \nu(L_T|_E)<r
 \label{altord:p:rank-drop}
\end{equation}
for every component of $D_N$ and every divisorial ramification
component of $\pi$ mapping onto a divisor of $S_0$. For the latter,
the inertia group of the Galois extension $\C(T)/\C(S_0)$ contains
a nonidentity element of $G_f$ fixing the divisor pointwise. These
fixed loci were included in $D_T$. Positive rational rescaling from
$J_T$ to $L_T$ leaves all these numerical dimensions unchanged.
The use of the actual image, before Stein factorization, is exactly
the function-field hypothesis needed in the finite-inertia part of
Theorem~\ref{thm:rankloss}.

\medskip\noindent
\textbf{Removal of the boundary and ramification.}
Write the effective Jacobian divisor as
\[
  R_\pi=K_T-\pi^*K_{S_0}
       =R_{\mathrm{div}}+R_{\mathrm{exc}},
\]
where $R_{\mathrm{div}}$ is supported on components mapping onto
divisors of $S_0$, and each component of $R_{\mathrm{exc}}$ has image
of codimension at least two.  In the volume
Lemma~\ref{altord:e:volume-lemma}, take
\[
  J=L_T,\qquad
  D=D_N+R_{\mathrm{div}},\qquad
  A=K_T-R_{\mathrm{div}}
   =\pi^*K_{S_0}+R_{\mathrm{exc}}.
\]
All components of $D$ satisfy the rank inequality above, and
$A+D=K_T+D_N$ is big by \eqref{altord:p:log-big}.  The volume lemma gives
\[
  A+jL_T
  \sim_{\Q}\pi^*(K_{S_0}+jL_0)+R_{\mathrm{exc}}
  \quad\text{big for }j\gg0.
\]
Lemma~\ref{altord:e:exceptional-pullback} removes the effective divisors
exceptional over $S_0$ and descends bigness through the generically
finite map.  Therefore
\begin{equation}
  K_{S_0}+jL_0\quad\text{is big for every sufficiently large integer }j.
  \label{altord:p:adjoint-big}
\end{equation}

\medskip\noindent
\textbf{Connected fibres.}
Finally take the Stein factorization of $p_0$, and resolve its finite
normal intermediate variety to a smooth projective variety $S$.
Resolving the induced rational map from $Y$ gives, by further smooth
projective birational modification, a morphism $p:Y\to S$ with
connected fibres. Indeed, its Stein factor is finite birational over
the normal variety $S$, because $\C(S)$ is already algebraically closed
in $\C(Y)$, and is therefore $S$ itself.
If $\sigma:S\to S_0$ is the resulting generically
finite morphism, set $L=\sigma^*L_0$.  It is nef, and
\eqref{altord:p:line-descent} pulls back to $M\sim_{\Q}p^*L$.
Since both $S$ and $S_0$ are smooth, the Jacobian divisor of $\sigma$
is effective.  Thus
\[
  K_S+jL
  =\sigma^*(K_{S_0}+jL_0)+R_\sigma
\]
is big for $j\gg0$ by \eqref{altord:p:adjoint-big} and preservation of
bigness under generically finite pullback.  This proves the theorem.
\end{proof}

\section{Canonical bundle data and the least-index root cover}\label{sec:canonical}
\label{altord:sec:canonical}

We return to the original algebraic fibre space. The goal of this
section is to construct the rational line to which
Theorem~\ref{thm:adjoint} applies and to retain the exact section
comparison with the source. The generic fibre of a relative Iitaka
fibration has Kodaira dimension zero; its least-index canonical
root cover provides the entire highest Hodge line.

We work over \(\C\). All varieties in this section are projective.
A rational pair $(Y,B)$ is klt if, on a log resolution, the boundary
defined by crepant pullback of $K_Y+B$ has every coefficient less
than one; it is log canonical if those coefficients are at most
one. We allow negative coefficients when explicitly discussing a
sub-pair. A b-divisor records compatible divisor traces on all
birational models of the base. Saying that its moduli part descends
to $Y$ means that every higher trace is the pullback of its trace $M$
on $Y$.
We distinguish a divisor exceptional over a birational model of the
total space from a divisor whose image on the base has codimension at
least two.

\subsection{The canonical bundle formula and its normalization}

We use the following parabolic form of the canonical bundle formula as
a deep theorem.

\begin{theorem}[Parabolic canonical bundle formula]
\label{altord:cb:parabolic}
Let \(g_0:X_0\to Y_0\) be an algebraic fibre space between smooth
projective varieties, with geometric generic fibre of Kodaira dimension
zero.  After birational modifications of its source and base there is a
commutative diagram
\[
\begin{array}{ccc}
 X'&\xrightarrow{\alpha}&X_0\\
 {\scriptstyle g}\downarrow&&\downarrow{\scriptstyle g_0}\\
 Y&\xrightarrow{\beta}&Y_0
\end{array}
\]
with \(X'\) and \(Y\) smooth, and rational divisors \(B,M,R\) such that
\begin{equation}
\label{altord:cb:formula}
 K_{X'}\sim_{\Q}g^*(K_Y+B+M)+R.
\end{equation}
Here \(B\geq0\), its support has simple normal crossings, \((Y,B)\) is
klt, and \(M\) is nef.  Moreover,
\begin{equation}
\label{altord:cb:residual}
 g_*\mathcal O_{X'}(\lfloor kR^+\rfloor)=\mathcal O_Y
 \quad(k\in\Z_{>0}),
 \qquad
 \operatorname{codim}_Y g(\Supp R^-)\geq2,
\end{equation}
and \(R^-\) is \(\alpha\)-exceptional.
The divisor \(M\) is the trace of the moduli b-divisor, which can be
assumed to descend to \(Y\).  On a sufficiently high generically finite
base change admitting semistable reduction, its pullback is
\(\Q\)-linearly equivalent to a positive rational multiple of the
Schmid extension of the distinguished canonical Hodge eigensheaf of
the canonical root cover.
\end{theorem}

The residual divisor and discriminant assertions, with the normalization
by the least nonzero pluricanonical index, are
\cite[\S4.4 and Theorem~4.5(i)--(v)]{FujinoMori}.
For original boundary zero, the logarithmic semistable part agrees
with the semistable part of Section~2 by
\cite[Proposition~4.7]{FujinoMori}.
The moduli interpretation follows from the discriminant-threshold
definition in \cite[Definition~4.3]{FujinoMori} and the following Ambro
results.  The b-divisor and Hodge assertions use
\cite[Theorems~0.2 and~4.4(1),(3), Lemma~5.2(4)--(5), and
Propositions~5.4--5.5]{AmbroBoundary};
see also \cite[Proposition~3.1]{AmbroModuli}.
These statements allow the auxiliary sub-boundary to have negative
coefficients.

We record how the hypotheses of these results fit together.
In the notation of \cite[\S4.4]{FujinoMori}, applied with original boundary
zero, the crepant boundary upstairs is the negative of the effective
relative canonical divisor over the smooth original source.
Thus the choice \(\Xi=0\) satisfies the hypotheses of their
Theorem~4.5, including those of part~(v).
The coefficients of our \(B\) are their \(s_P/b\); part~(v) places them
in \([0,1)\).  Their residual divisor is \(bR\).
We verify the all-degree direct-image condition from the normalization
below, so that it remains available after the final change of model.
For the exceptionality assertion, flatten before resolving the main
transform, as in \cite[\S4.4]{FujinoMori}.  The inverse image of a base
subset of codimension at least two has codimension at least two in
the flat main transform; finite normalization preserves this bound.
A prime divisor of its resolution lying over that subset is therefore
exceptional over the normalized main transform.  This transform maps
birationally to \(X_0\), so the prime is also \(\alpha\)-exceptional.
Apply this to every component of \(R^-\).  Its small image on \(Y\)
alone would not imply exceptionality over \(X_0\).

Here is also the relevant normalization of the associated sub-pair.
On the geometric generic fibre \(C\), let \(b\) be the least positive
integer with \(H^0(C,bK_C)\ne0\), and let \(D_C\) be the divisor of its
unique section up to scalar.  The horizontal restriction is
\[
 R_C=\frac1bD_C,\qquad
 \Delta_C=-R_C.
\]
The least positive integer making \(K_C+\Delta_C\) principal is
also \(b\). Indeed, \(bK_C-D_C\) is principal. Conversely, if
\(m(K_C-D_C/b)\) is principal for a positive integer \(m\), then
\(mD_C/b\) is an integral effective divisor linearly equivalent to
\(mK_C\). Thus \(H^0(C,mK_C)\ne0\), which forces \(m\geq b\).
This identifies the least-index normalization in the sub-pair
construction with the pluricanonical index used for the root cover.
A rational relative pluricanonical section first defines a divisor
\(R_0\) with this restriction.  For a prime divisor \(P\subset Y\),
subtract from \(R_0\) the pullback of \(P\) multiplied by
\[
 a_P=\min_{E\mapsto P}
       \frac{\operatorname{coeff}_E(R_0)}
            {\operatorname{mult}_E(g^*P)}.
\]
Only finitely many \(a_P\) are nonzero.  The resulting residual divisor
is nonnegative over the generic point of \(P\), with coefficient zero
on at least one component over \(P\).

This normalization also gives the direct-image condition in every
positive degree. Let \(V\subset Y\) be a nonempty open subset, let
\(k>0\) be an integer, and let \(\varphi\) represent a section of
\(\mathcal O_{X'}(\lfloor kR^+\rfloor)\) over \(g^{-1}(V)\).
On the geometric generic fibre \(C\), its poles are bounded by
\(\lfloor kD_C/b\rfloor\), hence by an integral multiple of
\(D_C\). Since \(D_C\sim bK_C\) and \(\kappa(C)=0\), every such
function is constant on \(C\). Since \(g_*\mathcal O_{X'}=\mathcal O_Y\),
the field \(\C(Y)\) is algebraically closed in \(\C(X')\).
Thus \(\varphi=g^*\psi\) for a rational function \(\psi\) on
\(V\). For each prime divisor \(P\) meeting \(V\),
choose a component \(E\) over \(P\) on which the normalized
residual coefficient is zero. Regularity of \(\varphi\) along
\(E\) gives
\[
 0\leq\operatorname{ord}_E(g^*\psi)
   =\operatorname{mult}_E(g^*P)\operatorname{ord}_P(\psi).
\]
Thus \(\psi\) has no codimension-one pole on \(V\), and normality
of \(Y\) makes it regular. Pullbacks of regular functions give the
reverse inclusion, proving
\[
 g_*\mathcal O_{X'}(\lfloor kR^+\rfloor)=\mathcal O_Y
 \qquad(k\in\Z_{>0}).
\]

For the sub-pair \((X',-R)\) its discriminant coefficient is \(1-t_P\),
where
\[
 t_P=\sup\{t\in\R:
       (X',-R+t\,g^*P)\text{ is sub-lc over the generic point of }P\}.
\]
Nonnegativity of \(R\) there implies \(t_P>0\); the component with
residual coefficient zero gives \(t_P\leq1\).
This explains the effective boundary in this normalization.

The rank-one hypothesis for this sub-pair holds even though
\(\Delta_C\) is negative.  Indeed, \(\kappa(C,R_C)=0\).
On a log resolution of \((C,-R_C)\), a section of the rounded
discrepancy sheaf pushes down to a rational function on \(C\) whose
poles are bounded by a sufficiently large integral multiple of
\(D_C\).  The space of such functions has dimension one, while the
constant function is allowed.  Consequently
\[
 \rk g_*\mathcal O_{X'}\bigl(\lceil\mathbf A(X',-R)\rceil\bigr)=1,
\]
where the expression is interpreted by pushing down from a log
resolution.  This is the rank condition in
\cite[introductory assumptions (1)--(3)]{AmbroBoundary}.

After choosing a model on which the moduli b-divisor descends, take
a common higher base model satisfying the flattening and
normal-crossing requirements, with boundary containing the total
transform of the original normalized discriminant support and all
exceptional base divisors. Crepantly pull back the auxiliary sub-pair
to a smooth resolved main transform, taking a log resolution that
includes its residual support and the pullbacks of the base boundary.
Apply the same minimum-coefficient normalization to this residual divisor. Its
horizontal restriction is again \(D_C/b\) for the unique
least-index pluricanonical divisor on the new geometric generic
fibre: for the induced birational morphism \(\lambda:C'\to C\),
crepant pullback gives
\[
 \lambda^*(D_C/b)+K_{C'/C}=D_{C'}/b.
\]
The preceding argument therefore restores the all-degree
direct-image condition on this final model. The normalized residual
is nonnegative over every codimension-one base point; its negative
part has base image of codimension at least two and is exceptional
over the original source by the flattening argument above. Over the
strict transform of an original base prime, its zero-coefficient
component survives and the crepant residual stays nonnegative, so
renormalization makes no further change there. Any new discriminant
support is therefore exceptional and lies in the chosen
simple-normal-crossing boundary. Its coefficients are in \([0,1)\)
by the same threshold calculation, so the discriminant is effective
and klt.

This renormalization does not change the moduli b-divisor. Indeed,
changing a sub-boundary from \(\Delta\) to \(\Delta+g^*A\) changes
the divisor pulled back in \(K_{X'}+\Delta\) by \(A\), while the
discriminant changes by \(A\) as well, by its threshold definition.
Their difference is unchanged, and the same calculation holds on
every higher base model. Crepant birational pullback likewise
preserves the induced b-divisor. Thus the final normalized residual
retains the descended nef moduli part together with the effective
klt discriminant and the residual properties needed for the section
comparisons below.
The parabolic pullback statement and independence of the chosen root
trivialization are also given directly in
\cite[Definition-Proposition~7.1]{AmbroBoundary}.

\begin{remark}
\label{altord:cb:no-abundance}
No abundance assertion for \(M\) is part of the input.
The stronger conclusion of
\cite[Theorem~3.3]{AmbroModuli} assumes effectivity of the boundary on
the geometric generic fibre, an assumption which need not hold for
\(-R_C\).  The b-nefness and Hodge identification used above instead
come from the results of \cite{AmbroBoundary} just cited.
In particular, no good minimal model of \(C\) is required.
For the Hodge-theoretic formulation with sub-pairs, see also
\cite[Definition~6.18, Construction-Definition~6.23,
Remark~6.25, and Theorem~6.28]{BFMT}.
\end{remark}

\subsection{The least-index root cover and its rational variation}

\begin{lemma}[Canonical root cover]
\label{altord:cb:root-cover}
Let \(C\) be a smooth connected projective variety with
\(\kappa(C)=0\), and let \(b>0\) be the least integer for which
\(H^0(C,bK_C)\ne0\).  Choose a nonzero section \(s\) of \(bK_C\) and a
nonzero rational canonical form \(\omega\).
Let \(T\) be the normalization of \(C\) in the field obtained by
adjoining a root
\[
 \xi^b=\frac{s}{\omega^b}.
\]
Then \(T\to C\) has degree \(b\), and every smooth projective
resolution \(W\) of \(T\) satisfies
\[
 \kappa(W)=0,\qquad h^0(W,K_W)=1.
\]
The unique canonical line is spanned by the form
\(\xi h^*\omega\), where \(h:W\to C\) is the induced morphism.
\end{lemma}

\begin{proof}
Put \(n=\dim C\) and \(D_C=\operatorname{div}(s)\).
The case \(n=0\) is immediate.  For \(n>0\), first note that every
nonzero pluricanonical space on \(C\) has dimension one:
two independent sections in the same degree would give a nonconstant
rational map and hence positive Kodaira dimension.

The field \(\C(C)\) contains all roots of unity.
If \(s/\omega^b=u^p\) for a prime \(p\mid b\), the rational
\((b/p)\)-canonical form \(u\omega^{b/p}\) has \(p\)-th power \(s\).
At every prime divisor its order is \(p^{-1}\) times the nonnegative
order of \(s\), so this form has no poles.  A rational section of a
line bundle on the smooth variety \(C\) with no divisorial poles is
regular.  This contradicts the minimality of \(b\).
The Kummer irreducibility criterion now gives a field extension of
degree \(b\).  Denote the normalization map by \(\pi:T\to C\), and a
smooth projective resolution by \(\rho:W\to T\), so \(h=\pi\rho\).

The differential of \(h\) gives a nonzero morphism
\(h^*K_C\to K_W\) at the generic point, and a regular morphism of
line bundles everywhere.  Use this morphism and its tensor powers
when pulling back canonical forms.  The rational canonical form
\(\eta=\xi h^*\omega\) then satisfies
\[
 \eta^{\otimes b}=h^*s.
\]
The right side is a regular section of \(bK_W\).
Thus \(b\,\operatorname{ord}_E(\eta)\geq0\) for every prime divisor
\(E\subset W\).  Normality of \(W\) again implies that \(\eta\) is
regular, and therefore \(h^0(W,K_W)\geq1\).

We next bound the canonical divisor from above.
Choose \(K_C=\operatorname{div}(\omega)\), and on \(T\) use the
canonical Weil divisor determined by the differential pullback of
\(\omega\).  At a prime divisor \(P\subset C\) outside \(\Supp D_C\),
the order of \(s/\omega^b\) is divisible by \(b\).  Locally over the
discrete valuation ring at \(P\), after dividing the root by a
power of a uniformizer, the defining equation is a \(b\)-th root of
a unit.  It is unramified, since the characteristic is zero.
Consequently all ramification divisors lie over \(\Supp D_C\).

If \(Q\subset T\) lies over \(P\subset\Supp D_C\), write
\(e=e(Q/P)\) and \(d=\operatorname{ord}_P(D_C)\geq1\).
Tameness gives the ramification coefficient \(e-1\), and
\[
 e-1\leq ed=\operatorname{ord}_Q(\pi^*D_C).
\]
The canonical divisor formula for finite maps, interpreted at
codimension-one points, therefore gives
\[
 K_T\leq\pi^*(K_C+D_C).
\]
On \(W\) the same inequality holds away from the
\(\rho\)-exceptional divisors.  Hence there is an effective integral
\(\rho\)-exceptional divisor \(E\) such that
\begin{equation}
\label{altord:cb:cover-upper-bound}
 K_W\leq h^*(K_C+D_C)+E.
\end{equation}
This argument does not require a pullback of the Weil divisor \(K_T\).

Set \(A=K_C+D_C\).  Normality of \(T\) implies
\(\rho_*\mathcal O_W(mE)=\mathcal O_T\) for every \(m\geq0\):
a rational function allowed poles only on exceptional divisors is
regular in codimension one on \(T\), hence regular on \(T\).
Projection formula and \eqref{altord:cb:cover-upper-bound} imply
\[
 \kappa(W,K_W)
 \leq \kappa(W,h^*A+E)
 =\kappa(T,\pi^*A)
 =\kappa(C,A).
\]
The last equality is invariance of Iitaka dimension under a finite
surjective pullback.  In this situation it can be seen directly
from section rings: a homogeneous section upstairs satisfies its
monic characteristic polynomial over \(\C(C)\); its coefficients,
in degree \(jm\), are sections of \(jmA\), by the valuation bounds
for sections of \(m\pi^*A\).  Thus the upstairs graded ring is
integral over the downstairs ring, and their fraction fields have
the same transcendence degree.
Finally \(D_C\sim bK_C\), so
\[
 \kappa(C,A)=\kappa(C,(b+1)K_C)=0.
\]
Together with the nonzero form \(\eta\), this proves
\(\kappa(W)=0\), and then \(h^0(W,K_W)=1\).
\end{proof}

The underlying cover of Lemma~\ref{altord:cb:root-cover} is the
canonical root cover used in
\cite[Remark~2.6 and Theorem~3.1]{FujinoMori}.
Over a dense open subset \(U\subset Y\), a rational relative
pluricanonical section constructs this cover in families; a
projective resolution, followed by shrinking \(U\), gives a smooth
projective morphism \(h_U:\mathcal W_U\to U\).
Its geometric generic fibre has the invariants in
Lemma~\ref{altord:cb:root-cover}. Constancy of Hodge numbers in the smooth
projective family then gives geometric genus one on every fibre.
If their dimension is \(n\), then
\[
 \mathbb V=R^n(h_U)_*\Q,\qquad
 \rk F^n(\mathbb V\otimes\mathcal O_U)=1.
\]
The variation has an integral lattice modulo torsion and is
polarizable, using a relative ample class and the Lefschetz
decomposition.  Alternatively, its primitive middle cohomology
is polarizable and contains all of \(H^{n,0}\), since cup product
with a \((1,1)\)-class annihilates \(H^{n,0}\).
The cyclic deck transformation \(\xi\mapsto\zeta\xi\) acts on the
unique top form through \(\zeta\).  Thus the distinguished
canonical eigensheaf is the entire top Hodge line of a rational
polarizable variation, even when the character itself is not
defined over \(\Q\).  We do not replace the rational variation by
that character eigenspace.

After a smooth projective alteration \(\tau:T_Y\to Y\), resolution
of the boundary, and semistable reduction in codimension one, the
identification in Theorem~\ref{altord:cb:parabolic} reads
\begin{equation}
\label{altord:cb:hodge-realization}
 \tau^*M\sim_{\Q}aJ,\qquad a\in\Q_{>0},
\end{equation}
where \(J\) is the Schmid extension of the top Hodge line of
\(\tau^*\mathbb V\), with unipotent local monodromy.
Birational changes of the resolved root cover do not change its
canonical direct image: on a common resolution, differential
pullback identifies its canonical forms.
The extension comparison is the one in
\cite[Theorem~4.4(1),(3), Lemma~5.2(5), and
Proposition~5.5]{AmbroBoundary}, rather than
an identification asserted only on \(U\).
Changing the chosen rational trivialization by a function from
the base gives the same identification after a further finite
base change extracting the corresponding root.

\subsection{Relative Iitaka reduction and section comparisons}

\begin{proposition}[Canonical bundle reduction]
\label{altord:cb:reduction}
Let \(f:X\to Z\) be an algebraic fibre space between smooth connected
projective varieties, and let its geometric generic fibre \(F\) have
\(d=\kappa(F)\geq0\).
There exist smooth projective varieties \(\widetilde X,Y\),
a birational morphism \(\mu:\widetilde X\to X\), and algebraic fibre
spaces
\[
 \widetilde X\xrightarrow{g}Y\xrightarrow{q}Z,\qquad
 qg=f\mu,
\]
such that \(\dim Y-\dim Z=d\) and the geometric generic fibre of
\(g\) has Kodaira dimension zero.
They admit the data \(B,M,R\) of Theorem~\ref{altord:cb:parabolic},
with \(R^-\) exceptional over \(X\).
The moduli divisor \(M\) has the realization
\eqref{altord:cb:hodge-realization} for a polarizable rational variation
whose top Hodge piece has rank one.
Writing
\[
 D=K_Y+B+M,
\]
one has
\begin{equation}
\label{altord:cb:two-comparisons}
 \kappa(X)\geq\kappa(Y,D),
 \qquad
 D|_{Y_{\bar\eta}}\text{ is big},
\end{equation}
where \(\bar\eta\) is a geometric generic point of \(Z\).
The same bigness holds on very general smooth fibres of \(q\).
\end{proposition}

\begin{proof}
Over \(\C(Z)\), take a sufficiently divisible pluricanonical map
of the generic fibre and its Iitaka fibration, including the
relative algebraic closure of the image field.
The Iitaka fibration theorem gives image dimension \(d\) and
Kodaira dimension zero for the geometric generic fibre of this
fibration.  Formation of its pluricanonical spaces commutes with
extension to an algebraic closure of \(\C(Z)\); thus these
statements can be checked geometrically.
Extend the rational map to projective models over \(Z\), resolve
source and target, and take the connected-fibre factorization.
This gives \(X_0\to Y_0\to Z\) with the required dimensions.
The second map has connected fibres as well: pushing
\(\mathcal O_{X_0}\) first to \(Y_0\) and then to \(Z\) gives
\(\mathcal O_Z\), because the original \(f\) and the birational
map to its smooth source have this property.

Apply Theorem~\ref{altord:cb:parabolic} to \(X_0\to Y_0\).
All its modifications can be performed over \(Z\).
Exceptionality over \(X_0\) implies exceptionality over \(X\).
Lemma~\ref{altord:cb:root-cover} and the subsequent family construction
give the required rational variation and its top line.
The two comparisons are proved in Lemma~\ref{altord:cb:comparison}
below.
\end{proof}

\begin{lemma}[Comparison of section growth]
\label{altord:cb:comparison}
For the factorization and canonical bundle data in
Proposition~\ref{altord:cb:reduction}, \eqref{altord:cb:two-comparisons} holds.
\end{lemma}

\begin{proof}
The formula gives
\[
 K_{\widetilde X}+R^-
 \sim_{\Q}g^*D+R^+.
\]
Write \(K_{\widetilde X}=\mu^*K_X+E_\mu\), using compatible
canonical divisors.  Since \(X\) and \(\widetilde X\) are smooth,
\(E_\mu\) is effective and \(\mu\)-exceptional.
The same is true of \(E_\mu+R^-\).
For sufficiently divisible \(m\), normality and projection
formula therefore give
\[
 H^0\bigl(\widetilde X,m(K_{\widetilde X}+R^-)\bigr)
 \simeq H^0(X,mK_X).
\]
Multiplication by the canonical section of \(mR^+\) injects
\(H^0(Y,mD)\) into the space on the left; here we also use
\(g_*\mathcal O_{\widetilde X}=\mathcal O_Y\).
This proves \(\kappa(Y,D)\leq\kappa(X)\).

Now make the flat base change \(\bar\eta\to Z\).
The varieties \(\widetilde X_{\bar\eta}\) and \(Y_{\bar\eta}\)
are smooth; the former is birational to the geometric generic
fibre \(F\) of \(f\).
Restrictions of canonical divisors agree, up to linear
equivalence, with the corresponding canonical divisors of these
fibres.  The restricted formula consequently implies
\[
 K_{\widetilde X_{\bar\eta}}
 \leq_{\Q}
 g_{\bar\eta}^*(D|_{Y_{\bar\eta}})
 +R^+|_{\widetilde X_{\bar\eta}}.
\]
Here \(\leq_{\Q}\) means that the difference is
\(\Q\)-linearly equivalent to an effective divisor.
Formation of a proper direct image commutes with this flat
base change, so \eqref{altord:cb:residual} yields, for sufficiently
divisible \(m\),
\[
 (g_{\bar\eta})_*
 \mathcal O_{\widetilde X_{\bar\eta}}
       (mR^+|_{\widetilde X_{\bar\eta}})
 =\mathcal O_{Y_{\bar\eta}}.
\]
Projection formula now gives
\begin{align*}
 h^0(\widetilde X_{\bar\eta},mK_{\widetilde X_{\bar\eta}})
 &\leq
 h^0\bigl(\widetilde X_{\bar\eta},
 m g_{\bar\eta}^*(D|_{Y_{\bar\eta}})
       +mR^+|_{\widetilde X_{\bar\eta}}\bigr)\\
 &=h^0(Y_{\bar\eta},mD|_{Y_{\bar\eta}}).
\end{align*}
It follows that
\[
 d=\kappa(F)
 \leq\kappa(Y_{\bar\eta},D|_{Y_{\bar\eta}})
 \leq\dim Y_{\bar\eta}=d.
\]
Thus \(D|_{Y_{\bar\eta}}\) is big, including the
zero-dimensional case.

For completeness, the bigness assertion spreads to a dense open
subset after restricting to the smooth locus of \(q\).
Fix a relatively ample divisor \(H\).
Bigness on the geometric generic fibre gives, for some
sufficiently divisible integer \(m\),
a nonzero section of \(mD-H\) there.  Flat scalar extension of
cohomology first gives such a section over \(\C(Z)\).
Shrinking the base, the section extends to the family and remains
nonzero on its fibres.  Hence \(mD|_{Y_z}\) is the sum of the
ample divisor \(H|_{Y_z}\) and an effective divisor for general
\(z\), proving the stated very-general-fibre assertion.
\end{proof}

\begin{lemma}[Further birational modifications]
\label{altord:cb:high-model}
Suppose \(Y,B,M,D\) satisfy the conclusions of
Proposition~\ref{altord:cb:reduction}.  Let \(\nu:Y'\to Y\) be a smooth
projective log resolution, and write
\[
 K_{Y'}+B^{\mathrm{cr}}=\nu^*(K_Y+B).
\]
Set \(B'=(B^{\mathrm{cr}})^+\), \(M'=\nu^*M\), and
\(D'=K_{Y'}+B'+M'\).
Then \(B'\) is effective, \((Y',B')\) is klt, \(M'\) is nef with
the same Hodge realization after pullback, and
\[
 D'=\nu^*D+(B^{\mathrm{cr}})^-.
\]
The additional divisor is effective and \(\nu\)-exceptional.
Consequently \(\kappa(Y',D')=\kappa(Y,D)\), and \(D'\) is big
on the geometric generic fibre over \(Z\).
\end{lemma}

\begin{proof}
The nonexceptional coefficients of \(B^{\mathrm{cr}}\) are those
of the effective divisor \(B\), so its negative part is
exceptional.  All coefficients of \(B^{\mathrm{cr}}\) are less
than one by the klt hypothesis.  Its positive part has simple
normal-crossing support and coefficients in \([0,1)\), proving
that \((Y',B')\) is klt.
The displayed identity follows directly from the definitions.
Exceptional effective twists do not change section spaces of a
pullback to a normal base, so they do not change the Iitaka
dimension.
On the geometric generic fibre over \(Z\), the induced map is
birational and the restricted identity expresses \(D'\) as the
pullback of a big divisor plus an effective divisor.
Nefness is preserved by pullback.  For the Hodge assertion, use
a common smooth alteration of \(Y'\) and the alteration in
\eqref{altord:cb:hodge-realization}; the unipotent Deligne and Schmid
extensions commute with the resulting morphisms of smooth
normal-crossing pairs.
\end{proof}

\begin{remark}\label{altord:cb:reduced-moduli-remark}
After Lemma~\ref{altord:cb:high-model}, \(B'\) need not be described as
the discriminant of the original sub-pair.  The properties used
later are precisely effectiveness and klt singularities, the
nef Hodge realization of \(M'\), and the two comparisons in
\eqref{altord:cb:two-comparisons}; all have been preserved.
In particular, since the big cone of \(Y_{\bar\eta}\) is open,
there is a rational number \(v\) with \(0<v<1\) such that
\[
 (K_Y+B+vM)|_{Y_{\bar\eta}}
 \quad\text{is big}.
\]
This remains true on very general fibres by the spreading
argument in Lemma~\ref{altord:cb:comparison}.
\end{remark}

\section{Ordinary subadditivity and extension of scalars}\label{sec:ordinary}\label{altord:sec:application}

We complete Theorem~\ref{thm:ordinary}. Section~\ref{sec:canonical}
has reduced the problem to an effective klt pair $(Y,B)$ with a nef
Hodge line $M$. The period theorem writes $M=p^*L$. The original
base map $q:Y\to Z$ and this new map $p:Y\to S$ need not factor
through one another. Their combined image will first give
nonvanishing on a period fibre. Weak positivity and one exact
interpolation identity then supply the missing total-space sections.

We first work over \(\C\).  Two established positivity results will be
used in precisely the following forms.

\begin{theorem}[Subadditivity with log-general-type fibre]
\label{altord:a:log-general-type}
Let \(u:V\to T\) be a surjective morphism with connected fibres, where
\(V\) is normal projective and \(T\) is smooth projective.  Let
\(\Delta\geq0\) be a \(\Q\)-divisor such that \((V,\Delta)\) is log
canonical.  If the log canonical divisor of the geometric generic
fibre \(G\) is big and \(\kappa(T)\geq0\), then
\[
 \kappa(V,K_V+\Delta)\geq \dim G+\kappa(T).
\]
\end{theorem}

This is the log-general-type-fibre theorem of Kov\'acs--Patakfalvi
\cite[Theorem~9.9]{KovacsPatakfalvi}, in the normal-total-space form
\cite[Theorem~2.11]{Hashizume}, with the auxiliary
divisor equal to \(K_T\).  It also follows from
\cite[Theorem~1.4(1)]{Hashizume}: a big divisor is abundant.  In
particular, the theorem does not require abundance of \(K_T\).

\begin{theorem}[Twisted weak positivity]\label{altord:a:weak-positivity}
Let \(u:V\to T\) be a surjective morphism from a normal projective
variety to a smooth projective variety, and let
\((V,\Delta)\) be an effective log canonical \(\Q\)-pair.  If
\(k\) is a positive integer such that \(k(K_V+\Delta)\) is Cartier,
then
\[
 u_*\mathcal O_V\bigl(nk(K_{V/T}+\Delta)\bigr)
\]
is weakly positive for every positive integer \(n\).
\end{theorem}

This is the projective case of \cite[Theorem~1.1]{FujinoWeakPositivity}.
The meaning needed here is that, for a torsion-free weakly positive
sheaf \(\mathcal E\), every ample Cartier divisor \(A\), and every
positive integer \(\alpha\), there is a positive integer \(\beta\)
such that
\[
 \bigl(\Sym^{\alpha\beta}\mathcal E\bigr)^{**}
       \otimes\mathcal O_T(\beta A)
\]
is generated by global sections at the generic point; see
\cite[Definitions~7.1--7.2]{FujinoWeakPositivity}.

\subsection{Nonvanishing on fibres of the period projection}

\begin{lemma}\label{altord:a:period-fibre}
Let \(p:Y\to S\) and \(q:Y\to Z\) be algebraic fibre spaces between
smooth projective varieties.  Let \((Y,B)\) be an effective klt
\(\Q\)-pair and let \(M\sim_{\Q}p^*L\) for a \(\Q\)-divisor \(L\) on
\(S\).  Suppose that \(\kappa(Z)\geq0\) and
\[
 D:=K_Y+B+M
\]
is big on the geometric generic fibre of \(q\).
Then the geometric generic fibre of \(p\) has nonnegative log
Kodaira dimension.
\end{lemma}

\begin{proof}
Take the Stein factorization of the map from \(Y\) to its image in
\(S\times Z\).  Resolve its intermediate normal variety and resolve
the induced rational map from \(Y\), choosing the modification of \(Y\)
also to be a log resolution of \((Y,B)\). We obtain morphisms
\[
 Y\xrightarrow{r}W,\qquad t:W\to S,\qquad h:W\to Z,
 \qquad p=t\circ r,\quad q=h\circ r,
\]
where \(Y,W\) are smooth projective, \(r\) has connected fibres, and
\(W\to S\times Z\) is generically finite onto its image.
We continue to write \(Y,B,M,D\) for the modified data.
This replacement is harmless: for such a log resolution
\(\mu:Y'\to Y\), take the positive part \(B'\) of the crepant
boundary.  Then
\[
 K_{Y'}+B'=\mu^*(K_Y+B)+E,\qquad E\geq0
\]
with \(E\) exceptional, and \((Y',B')\) is klt.
The same identity on the general \(p\)-fibres preserves their log
Kodaira dimensions.  Pullback of \(M\) and the effective correction
preserve the required bigness on the geometric generic \(q\)-fibre.

Since \(p_*\mathcal O_Y=\mathcal O_S\) and
\(r_*\mathcal O_Y=\mathcal O_W\), the fibres of \(t\) are connected.
For general \(s\in S\), the fibre \(W_s\) is smooth and connected,
and \(h_s:W_s\to Z\) is generically finite onto its image.  Write
\[
 a=\dim W_s,\qquad m=\dim Z.
\]
The generically finite map onto its image gives $a\leq m$.
We claim that \(\kappa(W_s)\geq0\).  The assertion is immediate for
\(a=0\); assume \(a>0\).  Along a smooth fibre of \(t\), the cotangent
sequence is
\[
 0\longrightarrow \Omega^1_{S,s}\otimes\mathcal O_{W_s}
 \longrightarrow \Omega^1_W|_{W_s}
 \longrightarrow \Omega^1_{W_s}\longrightarrow0.
\]
Since the last term has rank \(a\), the entire \(m\)-th exterior
power lies in the filtration term with at least \(m-a\) base
factors.  Its quotient by the next term gives, after pulling back
top forms on \(Z\), a morphism
\begin{equation}\label{altord:a:cotangent-map}
 h_s^*\omega_Z\longrightarrow
 \bigwedge\nolimits^{m-a}\Omega^1_{S,s}\otimes\omega_{W_s}.
\end{equation}
At a general point of a general \(W_s\), the differential of \(h\)
has rank \(m\), and its restriction to the tangent space of \(W_s\)
has rank \(a\).  Thus \eqref{altord:a:cotangent-map} is nonzero.  A suitable
linear functional on the constant exterior-power factor yields a
nonzero morphism \(h_s^*\omega_Z\to\omega_{W_s}\).

Choose \(n>0\) and \(0\ne\tau\in H^0(Z,nK_Z)\).
For general \(s\), the image \(h_s(W_s)\) is not contained in the
zero divisor of \(\tau\): otherwise dominance of \(h\) would fail.
The \(n\)-th tensor power of the preceding morphism therefore sends
\(h_s^*\tau\) to a nonzero section of \(nK_{W_s}\).  This proves the
claim.

Let \(G\) be a general fibre of \(r\).  For general \(z\in Z\), such
fibres form a moving family covering \(Y_z\).  A big divisor
restricts to a big divisor on a general member of this family:
restrict an ample-plus-effective decomposition, choosing a member
not contained in the support of the effective part.  Consequently
\(D|_G\) is big.  Since \(G\) maps to a point of \(S\),
\(M|_G\sim_{\Q}0\), and adjunction gives
\[
 D|_G\sim_{\Q}K_G+B|_G.
\]
For general \(s\), the morphism \(r_s:Y_s\to W_s\) has connected
fibres, \((Y_s,B|_{Y_s})\) is klt, and its general fibre has big log
canonical divisor.  Theorem~\ref{altord:a:log-general-type} gives
\[
 \kappa(Y_s,K_{Y_s}+B|_{Y_s})\geq
 \dim G+\kappa(W_s)\geq0.
\]
Finally choose \(s\) outside the countably many proper jumping
loci for the spaces of sections of all divisible log
pluricanonical bundles on the smooth \(p\)-fibres.  Generic base
change identifies their dimensions with those on the geometric
generic fibre.  Thus the last nonvanishing holds on that fibre as
well.
\end{proof}

\subsection{Effectivity from weak positivity}

\begin{lemma}\label{altord:a:effectivity}
Let \(p:Y\to S\) be an algebraic fibre space between smooth
projective varieties, and let \((Y,B)\) be an effective klt
\(\Q\)-pair.  Suppose that the geometric generic fibre has
nonnegative log Kodaira dimension.  For every big \(\Q\)-divisor
\(P\) on \(S\), the divisor
\[
 K_{Y/S}+B+p^*P
\]
is \(\Q\)-linearly equivalent to an effective divisor.
\end{lemma}

\begin{proof}
Choose a dense open over which \(p\) is flat. Flatten \(p\) by a
projective birational modification of the base
and then resolve the base, retaining flatness of the main
transform by base change; see \cite[Lemma~38.31.1, Tag~081R]{StacksFlatten}.
Resolve the resulting main transform \(V\to S_1\), choosing the
projective resolution so that its composite \(\mu:Y_1\to Y\) is also
a log resolution of \((Y,B)\). Thus
\[
 \sigma:S_1\to S,\qquad
 \mu:Y_1\to Y,\qquad
 p_1:Y_1\to S_1
\]
are compatible morphisms, \(S_1,Y_1\) are smooth projective,
\(\sigma,\mu\) are birational, and \(Y_1\to V\) is a resolution
with \(V\to S_1\) flat.
Any prime divisor of \(Y_1\) whose image in \(S_1\) has codimension
at least two is \(\mu\)-exceptional.  Indeed, flatness implies that
the inverse image of such a base subset has codimension at least
two in \(V\); hence the divisor is exceptional over \(V\), and
its image in \(Y\) also has codimension at least two.

Take the positive part \(B_1\) of the crepant boundary, so that
\((Y_1,B_1)\) is effective klt and
\begin{equation}\label{altord:a:source-correction}
 K_{Y_1}+B_1=\mu^*(K_Y+B)+E,\qquad
 E\geq0,\quad E\ \text{\(\mu\)-exceptional}.
\end{equation}
Generic-fibre nonvanishing and field-extension base change allow
us to choose a sufficiently divisible \(k>0\) for which
\[
 \mathcal E=
 p_{1*}\mathcal O_{Y_1}\bigl(k(K_{Y_1/S_1}+B_1)\bigr)
\]
has positive rank.  Theorem~\ref{altord:a:weak-positivity} applies to
\(\mathcal E\).

Fix an ample Cartier divisor \(A\) on \(S_1\), and choose a
positive integer \(\alpha\) so large that
\[
 \sigma^*P-\frac{1}{k\alpha}A
\]
is big.  For some \(\beta>0\), weak positivity makes
\[
 \bigl(\Sym^{\alpha\beta}\mathcal E\bigr)^{**}
       \otimes\mathcal O_{S_1}(\beta A)
\]
generically generated.  Set \(\ell=\alpha\beta\).
Let \(U\subset S_1\) be the locally free locus of the torsion-free
sheaf \(\mathcal E\); its complement has codimension at least two.
On all of \(U\), multiplication of sections gives
\[
 \Sym^\ell\mathcal E|_U\longrightarrow
 p_{1*}\mathcal O_{Y_1}
       \bigl(k\ell(K_{Y_1/S_1}+B_1)\bigr)|_U.
\]
This map is nonzero at the generic point: a nonzero section on the
integral generic fibre has nonzero powers.  Generic generation
therefore supplies a global section of the twisted reflexive
symmetric power whose product is nonzero.  Its image is a regular
section over \(p_1^{-1}(U)\) of
\[
 k\ell(K_{Y_1/S_1}+B_1)+\beta p_1^*A.
\]
Here \(U\) is the locally free locus, not merely a smaller open
where the chosen sections generate.

The big divisor
\[
 k\ell\sigma^*P-\beta A
   =k\ell\left(\sigma^*P-\frac{1}{k\alpha}A\right)
\]
has an effective sufficiently divisible multiple.  Taking the
same power of the preceding section and multiplying by the
pullback of a nonzero section of that multiple gives a nonzero
section over \(p_1^{-1}(U)\) of a multiple of
\[
 D_1:=K_{Y_1/S_1}+B_1+p_1^*\sigma^*P.
\]
As a rational section on \(Y_1\), it has poles only above
\(S_1\setminus U\).  Thus there are \(n>0\), a rational function
\(\varphi\ne0\), and an effective \(\mu\)-exceptional divisor \(T\)
such that
\begin{equation}\label{altord:a:section-with-poles}
 \operatorname{div}(\varphi)+nD_1+T\geq0.
\end{equation}

Choose compatible canonical divisors and put
\[
 D_0=K_{Y/S}+B+p^*P,\qquad
 R=K_{S_1}-\sigma^*K_S\geq0.
\]
Equation~\eqref{altord:a:source-correction} gives the exact identity
\[
 D_1=\mu^*D_0+E-p_1^*R.
\]
Push \eqref{altord:a:section-with-poles} forward to the smooth variety
\(Y\).  Since \(\mu_*E=\mu_*T=0\), we obtain
\[
 \operatorname{div}(\varphi)+nD_0
      \geq n\mu_*p_1^*R\geq0.
\]
This proves rational linear effectivity of \(D_0\).
\end{proof}

\subsection{Completion over the complex numbers}

\begin{proof}[Proof of Theorem~\ref{altord:ordinary-subadditivity} over \(\C\)]
If either \(\kappa(F)\) or \(\kappa(Z)\) is \(-\infty\), the
assertion follows from the convention in the statement.
Assume therefore that
\[
 d:=\kappa(F)\geq0,\qquad \kappa(Z)\geq0.
\]
Proposition~\ref{altord:cb:reduction} and Lemma~\ref{altord:cb:comparison}
provide algebraic fibre spaces
\[
 \widetilde X\longrightarrow Y\xrightarrow{q}Z
\]
and an effective klt pair \((Y,B)\), together with a nef
\(\Q\)-divisor \(M\) having the required Hodge realization.
Writing \(D=K_Y+B+M\), their conclusions include
\begin{equation}\label{altord:a:reduction}
 \dim(Y/Z)=d,\qquad
 \kappa(X)\geq\kappa(Y,D),\qquad
 D|_{Y_{\bar\eta}}\ \text{big},
\end{equation}
where \(Y_{\bar\eta}\) is the geometric generic \(q\)-fibre.

Apply adjoint positivity, Theorem~\ref{altord:p:observation},
and the compatible modifications of Lemma~\ref{altord:cb:high-model}.
We obtain an algebraic fibre space \(p:Y\to S\) with \(S\) smooth
projective and
\[
 M\sim_{\Q}p^*L,\qquad L\ \text{nef}.
\]
If \(\dim S=0\), then \(M\sim_{\Q}0\).
Theorem~\ref{altord:a:log-general-type}, applied directly to \(q\), and
\eqref{altord:a:reduction} give the result.
We may consequently suppose \(\dim S>0\).  Adjoint positivity also gives
\[
 K_S+jL\ \text{big for all sufficiently large integers }j.
\]

Openness of the big cone on \(Y_{\bar\eta}\) permits a rational
number \(v\) with \(0<v<1\) such that
\begin{equation}\label{altord:a:reduced-moduli}
 (K_Y+B+vM)|_{Y_{\bar\eta}}\ \text{is big}.
\end{equation}
Lemma~\ref{altord:a:period-fibre} shows that the geometric generic
\(p\)-fibre has nonnegative log Kodaira dimension.
Choose a sufficiently large integer \(j>1\) and an ample
\(\Q\)-divisor \(H\) on \(S\) such that
\[
 P_0:=K_S+jL-H
\]
is big.  Lemma~\ref{altord:a:effectivity} gives
\begin{equation}\label{altord:a:effective-term}
 E_1:=K_Y+B+jM-p^*H
      \sim_{\Q}K_{Y/S}+B+p^*P_0
      \ \sim_{\Q}\ \text{an effective divisor}.
\end{equation}

Choose
\[
 \lambda=\frac{1-v}{j-v}\in(0,1)\cap\Q,
 \qquad \lambda j+(1-\lambda)v=1.
\]
This balances the coefficients of \(M\); to cancel the corresponding
term \(-\lambda p^*H\), put
\[
 A_S=vL+\frac{\lambda}{1-\lambda}H,\qquad
 E_2=K_Y+B+p^*A_S.
\]
Then \(A_S\) is ample and
\begin{equation}\label{altord:a:interpolation}
 D\sim_{\Q}\lambda E_1+(1-\lambda)E_2.
\end{equation}
Choose a sufficiently divisible integer \(N>1\) such that
\(|Np^*A_S|\) is basepoint-free.  A general member
\(T_N\in|Np^*A_S|\) gives
\[
 \Gamma=\frac{1}{N}T_N\sim_{\Q}p^*A_S
\]
with \((Y,B+\Gamma)\) klt.  This follows from Bertini, or after
pulling the basepoint-free system to a log resolution of
\((Y,B)\) and choosing \(N\) sufficiently large.
Thus \(E_2\sim_{\Q}K_Y+B+\Gamma\) is the log canonical class of
an effective klt pair.  Its restriction to \(Y_{\bar\eta}\) is
the big divisor in \eqref{altord:a:reduced-moduli} plus the nef divisor
\(\lambda p^*H/(1-\lambda)\), and hence is big.
Theorem~\ref{altord:a:log-general-type} now yields
\[
 \kappa(Y,E_2)\geq d+\kappa(Z).
\]
In sufficiently divisible degrees, multiplication by powers of a
nonzero section in \eqref{altord:a:effective-term} embeds the corresponding
linear systems of \((1-\lambda)E_2\) into those of \(D\).
Positive rational scaling leaves Iitaka dimension unchanged, so
\[
 \kappa(X)\geq\kappa(Y,D)\geq\kappa(Y,E_2)
        \geq\kappa(F)+\kappa(Z),
\]
as required.
\end{proof}

\subsection{Extension to arbitrary characteristic-zero fields}

\begin{lemma}\label{altord:a:field-extension}
Let \(K\subset K'\) be a field extension and let \(V\) be a smooth
projective geometrically integral variety over \(K\).
The Kodaira dimension of \(V\) is unchanged by extending scalars
to \(K'\).  The same assertion holds for the geometric generic
fibre of an algebraic fibre space between smooth projective geometrically integral
varieties when the ground field is extended.
\end{lemma}

\begin{proof}
Flat base change for proper coherent cohomology gives, in every
positive degree $n$,
\[
 H^0(V,nK_V)\otimes_K K'
 \simeq H^0(V_{K'},nK_{V_{K'}}).
\]
The canonical bundle commutes with this smooth scalar extension.
Nonvanishing is therefore unchanged. When the space is nonzero,
choose a $K$-basis. The rational map defined by that basis and its
image base-change to the complete pluricanonical map over $K'$;
dimensions of finite-type schemes are unchanged by extension of
the ground field. Taking the supremum of these image dimensions
proves the first assertion, including the case $-\infty$.

For the fibre assertion, put \(E=K(Z)\) and
\(E'=K'(Z_{K'})\), and write \(F_K\) for the generic fibre over \(E\).
The generic fibre after ground-field extension is \(F_K\times_E E'\).
Choose an algebraic closure \(\Omega\) of \(E'\) and an algebraic
closure \(\overline E\subset\Omega\) of \(E\). The two geometric
generic fibres are \(F_K\times_E\overline E\) and
\(F_K\times_E\Omega\); the latter is a scalar extension of the
former. Their Kodaira dimensions agree by the first assertion.
\end{proof}

To complete the proof of Theorem~\ref{altord:ordinary-subadditivity}, let \(k\) be any
algebraically closed field of characteristic zero.
The varieties, their projective embeddings, and the morphism
descend to a finitely generated subfield \(K\subset k\) over
\(\Q\).  Smoothness, geometric integrality and surjectivity hold
for the descended data by faithful flatness.  Proper flat base
change and faithful flatness also preserve the equality
\(f_*\mathcal O_X=\mathcal O_Z\), hence the connected-fibre
condition.  Choose an embedding \(K\hookrightarrow\C\).
Lemma~\ref{altord:a:field-extension} identifies all three Kodaira
dimensions in the statement with their counterparts over
\(\C\), using the geometric generic fibre for \(F\).
The complex case just proved therefore establishes the theorem
over \(k\).

\section{An integral projective proof of adjoint positivity}
\label{p1logalt:section}

We give another proof of adjoint positivity for the whole highest Hodge
line of a rationally polarized integral pure variation.  The conclusion
is the same adjoint conclusion as Theorem~\ref{thm:adjoint} for the
rational parabolic extension itself. A positive rational multiple
has the same conclusion after rescaling the large adjoint parameter.
We first replace the
variation by an integral tensor construction, then descend it through
its full period map and subtract the boundary by comparing intersection
growth.

The tensor construction removes every monodromy element acting trivially
on the moving projective line.  This makes a transverse loop detect
strict loss of the line's numerical dimension at each remaining boundary
divisor.  The full period map supplies an algebraic base; the differential
of the map remembering only the highest line controls the boundary
subtraction.  These two ranks need not agree.

Let $P$ be a smooth connected projective complex variety, let $D$ be a reduced simple normal crossing
divisor, and put $P^\circ=P\setminus D$.  A \emph{top Hodge line} of a rationally
polarized integral pure variation $\mathbb V$ on $P^\circ$ means
\[
 F^p\mathbb V_{\cO}\quad\text{with}\quad
 F^{p+1}\mathbb V_{\cO}=0,\qquad
 \rank F^p\mathbb V_{\cO}=1.
\]
Here an integral variation has a local system of finite-rank free
$\mathbb Z$-modules, and the polarization is rational and flat.  The
weight and $p$ may be arbitrary integers.  The line itself is a complex
Hodge subbundle; no rationality assumption is made on it.  If $V$ is a flat reference
fiber, its \emph{line period map} is the equivariant holomorphic map
\[
 \ell:\widetilde{P^\circ}\longrightarrow\mathbb P(V),
 \qquad x\longmapsto F^p_x.
\]
Its generic differential rank does not depend on the reference fiber.

\begin{theorem}[Adjoint positivity]\label{p1logalt:positivity}
Let $P$ be a smooth connected projective complex variety and $D$ a
reduced simple normal crossing divisor.  Let $\mathbb V$ be a rationally
polarized integral pure variation on $P\setminus D$ whose highest
nonzero Hodge filtration piece is a line, and let
$L\in\operatorname{Pic}(P)\otimes\Q$ be its rational parabolic extension.  There are a projective birational
morphism $\pi:\widetilde P\to P$, with $\widetilde P$ smooth, a smooth
projective variety $Z$, a surjective morphism $h:\widetilde P\to Z$ with
connected fibers, and a nef rational divisor $A$ on $Z$ such that
\[
 \pi^*L\sim_{\Q}h^*A,
 \qquad K_Z+aA\ \text{is big for every sufficiently large rational }a.
\]
Here a divisor on a point is big.  Further smooth birational modifications
of the source are permitted by replacing $h$ with its composition with
the modification.
\end{theorem}

We first specify the extension and metric facts used in the proof.

\subsection{Extensions and numerical dimension}

The monodromy theorem makes the local monodromies quasi-unipotent.
For unipotent local monodromy, $L$ is the top filtration subbundle in the
Deligne--Schmid extension.  For quasi-unipotent monodromy, its rational
extension is characterized by the following rule: after an adapted cover
making local monodromy unipotent, its pullback is the unipotent extension,
as a rational line bundle.  This convention includes the parabolic
weights; it is not an arbitrary choice of a Deligne lattice.

\begin{lemma}[Extension facts]\label{p1logalt:extensions}
The rational extension has the following properties.
\begin{enumerate}[label=\textup{(\roman*)}]
\item It is nef and is compatible with positive tensor powers and
pullback by morphisms of smooth projective log pairs whose interiors map
to the original interior.
\item In the unipotent case, on each connected boundary stratum the local
limiting mixed Hodge structures determine a unique weight-graded piece
containing the top Hodge line.  These pieces glue to a polarizable pure
variation with rank-one top filtration, whose extended top line on the
closure of the stratum is the restriction of $L$.
\item The preceding restriction statement remains valid after an adapted
cover in the quasi-unipotent case.  In particular, numerical dimensions
of restrictions may be computed on a smooth adapted cover of the
stratum.
\end{enumerate}
\end{lemma}

\begin{proof}
In the unipotent case the canonical extension and relative monodromy
weight filtration are supplied by nilpotent-orbit theory
\cite{Schmid,CKS}; precise formulations for the top
line are given in \cite[\S 2.5.3, Lemma 2.16, Lemma 5.6]{BFMT}.  The
graded pieces are polarized by the primitive decomposition.

For a boundary stratum, take the weight filtration associated with an
interior point of its monodromy cone; it is independent of that choice.
Strictness and $\rank F^p=1$ select one nonzero highest graded piece.
The global object in (ii) is this pure graded variation.  To check its
gluing when the normal bundle is nontrivial, choose a local normal
coordinate $t$ and write $N$ for its monodromy logarithm.  Replacing $t$
by $t'=u(s)t$, with $u$ invertible along the stratum, changes the local
limiting filtration by
\[
 \exp\left(-\frac{\log u(s)}{2\pi\sqrt{-1}}N\right).
\]
Since $N$ lowers the monodromy weight filtration by two, this transition
acts identically on each weight graded.  For several normal coordinates
the same calculation uses the sum of the commuting monodromy logarithms,
which also acts trivially on their common weight graded.  The pure
graded data therefore glue.  Strictness identifies the rank-one top
filtration with its unique nonzero top graded piece.  At the remaining
corners, \cite[\S 2.5.3 and Lemma 2.16(3)]{BFMT} identifies the canonical
extensions on the entire smooth stratum closure, so the identification
of lines holds there as well.

For completeness, the rational convention is compatible with these
operations as follows.  In a crossing chart, take sufficiently divisible
roots of its boundary coordinates.  The resulting unipotent extension
is equivariant for the finite deck group.  A common positive power kills
the characters of the stabilizers on its line fibers and descends to a
line bundle in the original chart.  The descents agree on overlaps:
compare them on a common root cover and use the uniqueness of the
unipotent extension.  This defines a class in
$\operatorname{Pic}(P)\otimes\Q$ and proves compatibility with further
pullbacks and powers.  One can perform the comparison globally on a
smooth adapted alteration; auxiliary branch divisors cause no difficulty
because the original variation extends over them.  Nefness descends
under a proper surjective morphism, since every curve downstairs is
dominated by a curve upstairs.  Finally, nef numerical dimension is
unchanged by a generically finite pullback, so the same construction on
strata proves (iii).
\end{proof}

\begin{lemma}[Curvature and intersections]\label{p1logalt:rank}
The numerical dimension of the rational top Hodge line equals the
generic rank of its line period map:
\[
 \nu(L)=\rank(d\ell)_{\mathrm{gen}}.
\]
In the unipotent case, if $d=\dim P$, $\omega$ is a smooth K\"ahler form,
and $\alpha$ is the first Chern form of the Hodge metric on $L|_{P^\circ}$,
then, for $0\le k\le d$,
\begin{equation}\label{p1logalt:chern}
 \int_{P^\circ}\alpha^k\wedge\omega^{d-k}
 =c_1(L)^k\cdot[\omega]^{d-k}.
\end{equation}
The integrals are absolutely convergent.
\end{lemma}

\begin{proof}
The closed-test-form version of this calculation is proved in
Lemma~\ref{p1varhd:curvature}; that proof uses only the extension facts in
Lemma~\ref{p1logalt:extensions}, not any descent or rank-drop result.
Taking the test form to be $\omega^{d-k}$ gives
Equation~\eqref{p1logalt:chern}.  Griffiths' curvature formula gives
\begin{equation}\label{p1logalt:curvature-kernel}
                 \ker\alpha=\ker(d\ell).
\end{equation}
If $r$ is the generic differential rank, then $\alpha^{r+1}=0$ everywhere
on the open set, whereas $\alpha^r\wedge\omega^{d-r}$ is positive on a
nonempty open subset.  The intersection identity and nefness therefore
give $\nu(L)=r$.  For quasi-unipotent monodromy, both numerical dimension
and differential rank are unchanged by a smooth adapted generically
finite cover, and Lemma~\ref{p1logalt:extensions} identifies the pulled-back
rational line with its unipotent extension.
\end{proof}

\subsection{Eliminating the projective monodromy kernel}

For a variation on a connected smooth algebraic variety, we will use the
condition
\begin{equation}\label{p1logalt:faithful-line}
 \gamma\ell(x)=\ell(x)\text{ for every }x
 \quad\Longrightarrow\quad \gamma=\id
 \quad(\gamma\text{ in the monodromy image}).
\end{equation}
This is a condition on the image of the moving line, rather than on all
of the ambient projective space.

\begin{lemma}[An integral tensor replacement]\label{p1logalt:kernel}
Let $U$ be a smooth connected complex algebraic variety, and
let $\mathbb V$ be a rationally polarized integral pure variation on $U$
whose highest nonzero Hodge filtration piece is a line $L^\circ$.
There are an integer $b>0$ and a rationally polarized integral pure
variation $\mathbb V_1$ on $U$ satisfying
Condition~\eqref{p1logalt:faithful-line}, whose highest nonzero Hodge
filtration piece is $(L^\circ)^{\otimes b}$.  On any smooth projective
simple-normal-crossing compactification of $U$, the rational parabolic
highest line of $\mathbb V_1$ is the $b$th power of that of $\mathbb V$.
\end{lemma}

\begin{proof}
For a general algebraic $U$, choose a dense quasi-projective open
$U_0$.  The map $\pi_1(U_0)\to\pi_1(U)$ is surjective, since loops
can be moved off the complement.  Thus the monodromy decompositions
below are unchanged, and a flat projector of Hodge type $(0,0)$ on
$U_0$ has that type throughout $U$ by continuity.
Write $\Gamma$ for the original monodromy image.  By semisimplicity
\cite[Theorem~4.2.6 and the added-in-proof note]{DeligneHodgeII}, the complex local system is a sum of
monodromy-isotypic summands $V_j$. These are complex Hodge subvariations. To see this
precisely, the fixed-part theorem equips
$\operatorname{End}_{\Gamma}(V)$ with its Hodge algebra structure, and its
center has type $(0,0)$ by \cite[Corollary 4.2.8(i),(ii)]{DeligneHodgeII}.
Equivalently, the connected Hodge action fixes every primitive central
idempotent of this semisimple algebra.  The images of those idempotents
are exactly the $V_j$.

The top line belongs to one summand $V_{j_0}$.  Retain that summand and
all its rational conjugates.  Their sum is a rational subvariation:
the central idempotents are defined over a finite splitting field of the
rational commutant algebra, and the selected sum is Galois invariant.
It inherits the rational polarization and the lattice obtained by
intersecting with the original lattice.  Denote it by $V_0$.  None of
this construction requires the highest line, or its scalar character,
to be rational.

Suppose $\gamma\in\Gamma$ fixes the moving line projectively everywhere.
Its projective fixed locus is the finite union of the projectivizations
of its eigenspaces.  The connected line image lies in one of them, so
$\gamma$ is scalar on the linear span $S$ of that image.  Equivariance
makes $S$ a nonzero $\Gamma$-subrepresentation of $V_{j_0}$.  Write
$V_{j_0}=E\otimes M$, with $E$ irreducible and $\Gamma$ acting trivially
on $M$.  Since $S$ contains a copy of $E$, scalar action on $S$ implies
the same scalar action on $E$ and hence on all of $V_{j_0}$.  Rational
conjugation gives scalar action on every retained $V_j$.

It follows that $\gamma|_{V_0}$ preserves the Hodge filtration at any
point.  It also preserves the rational polarization and the lattice.
The stabilizer of a polarized Hodge structure in its real polarization
group is compact, since it preserves the positive definite Hodge metric.
Its intersection with the integral automorphism group is finite.
Therefore $\gamma|_{V_0}$ has finite order. For this particular
variation, one may take the exponent of that finite stabilizer at a
fixed base point. There is also a choice depending only on
$R=\rank V_0$.  Indeed, an eigenvalue of an integral finite-order matrix
of rank $R$ has order $q$ with $\varphi(q)\le R$, because its cyclotomic
minimal polynomial divides the characteristic polynomial.  The finite
set of such $q$ therefore gives the explicit permissible choice
\[
             b=\operatorname{lcm}\{q\geq1:\varphi(q)\leq R\}.
\]
Every scalar on every retained summand has $b$th power one.  This
uniform exponent is chosen before constructing the new variation.

Inside $V_0^{\otimes b}$ take
\[
 V_1=\bigoplus_{j\ \mathrm{retained}}V_j^{\otimes b}.
\]
The sum is taken inside $V_0^{\otimes b}$. The individual terms are
distinct direct summands in the tensor expansion, so the sum
is direct.  It is Galois invariant and hence rational, and each summand
is a Hodge subvariation.  Restricting the tensor polarization and
intersecting with the tensor lattice make it rationally polarized and
integral.  The summand
$V_{j_0}^{\otimes b}$ has top line $(L^\circ)^{\otimes b}$; all other
summands have strictly smaller highest Hodge index.  Thus this is the
unique top line of $V_1$.

Every element considered above acts trivially on $V_1$.  Conversely, if
an element of the new monodromy image fixes its moving top line, choose
a lift in $\Gamma$.  Injectivity of the Veronese map
$\mathbb P(V_{j_0})\to\mathbb P(V_{j_0}^{\otimes b})$ shows that the lift
fixes the original moving line.  It is therefore killed on $V_1$.
This proves Condition~\eqref{p1logalt:faithful-line}.  Compatibility of extensions
with tensor powers finishes the proof.
\end{proof}

We now descend the variation itself, using faithfulness to kill
monodromy along connected period fibres. Consequently we may
normalize the period image in the relative algebraic closure of its
function field in $\C(P)$ before estimating boundary ranks. The
finite-inertia contradiction in Section~\ref{sec:rankloss} required
the actual image; here the tensor replacement makes fibre monodromy
trivial once it fixes the moving line on an open set.

\begin{lemma}[Full-period descent]\label{p1logalt:descent}
After applying Lemma~\ref{p1logalt:kernel} and making smooth birational
modifications of the projective models, the resulting variation
descends over a dense open subset of a smooth projective variety $Z$.
Its descended full period map is generically immersive and satisfies
Condition~\eqref{p1logalt:faithful-line}.  If $A_1$ is its rational top Hodge
extension, then the original extension satisfies
\[
 b\pi^*L\sim_{\Q}h^*A_1
\]
for a morphism $h:\widetilde P\to Z$ with connected fibers.
\end{lemma}

\begin{proof}
Let $\mathcal D$ be the period domain for the chosen polarized lattice,
and choose an arithmetic group $\Gamma_{\mathrm{ar}}$ containing the
actual monodromy image $\Gamma$.  By \cite[Theorem 1.1]{BakkerBrunebarbeTsimerman}, the full
period map factors through an algebraic dominant morphism
$P^\circ\to Y_0$, where $Y_0$ is an irreducible quasi-projective variety
and $Y_0^{\mathrm{an}}\to\Gamma_{\mathrm{ar}}\backslash\mathcal D$ is
a closed analytic immersion.  Normalize $Y_0$ in the relative algebraic
closure of $\C(Y_0)$ in $\C(P)$.  This extension is finite, so the
normalization is finite over $Y_0$.  Its function field is relatively
algebraically closed in $\C(P)$; in characteristic zero this makes the
generic fiber geometrically connected.  Shrink
source and target so that the latter is smooth and the map is a smooth
topologically locally trivial fibration with connected fibers.  To
obtain this shrinking for the possibly nonproper source, stratify a
proper algebraic compactification of the map together with its source
boundary and restrict over the open base stratum.

A local lift of the full period map along a fiber takes values in one
$\Gamma_{\mathrm{ar}}$-orbit.  These orbits are discrete, including when
the quotient has finite stabilizers, so that lift is locally constant
in fiber directions.  Continuation around a fiber loop therefore fixes
the lifted full period.  Transport the loop to nearby fibers using the
local trivialization.  In flat coordinates its actual monodromy fixes
the top line on a nonempty analytic open subset of the source.  This
identity extends to its connected universal cover, and
Lemma~\ref{p1logalt:kernel} kills that monodromy.  The homotopy sequence of
the fibration now factors the integral representation through the
ordinary fundamental group of the base.  The filtration descends as well: on a
simply connected base chart the pulled-back flat bundle is trivial,
the filtration is constant on each connected fiber, and holomorphic
local sections of the smooth map show that the descended filtration is
holomorphic.  Its transversality and polarization follow by pullback.

The descended period map has the dimension of the original full period
image and is generically immersive.  Base loops lift to the source, so
the same argument and Condition~\eqref{p1logalt:faithful-line} show that the descended
variation retains that condition.  Shrinking introduces no change to
it: any monodromy element fixing the line on the smaller open fixes it
identically by analytic continuation.

Choose a smooth projective compactification of the base, resolve its
boundary, and resolve the rational map from $P$ to this compactification.
The resulting morphism $h:\widetilde P\to Z$ has connected generic
fiber and hence connected fibers, by Stein factorization and normality
of $Z$.  Enlarge the source boundary to contain the inverse image of
the chosen base boundary and resolve it.  The rational extension of
the original variation is unchanged where it already extended.
Lemma~\ref{p1logalt:extensions} identifies its $b$-th power with $h^*A_1$.
\end{proof}

\subsection{A strict rank drop at nontrivial monodromy}

\begin{lemma}[A transverse-loop criterion]\label{p4int:transverse-loop}
Let $U=\Delta\times\Delta^{d-1}$, with coordinates $(z,s)$, and let
$U^*=\Delta^*\times\Delta^{d-1}$.  Suppose that a holomorphic map
$\ell:\widetilde{U^*}\to\mathbb P(V)$ is equivariant for a linear
transformation $T\in\operatorname{GL}(V)$ under one positive turn around
$z=0$, and that $\rank(d\ell)\leq r$.  Suppose that a holomorphic map
$q:U\to\mathbb C^r$ has tangential differential of rank $r$ at a point
of $z=0$, and that on $U^*$ its lift satisfies
$\ker(d\ell)\subseteq\ker(dq)$.
Then $T$ fixes every line in the image of $\ell$ on a nonempty open
subset.  If $\ell$ is the restriction of a line period map on a connected
universal cover, $T$ fixes that entire line image.  This includes $r=0$,
when $q$ is the map to a point.
\end{lemma}

\begin{proof}
After shrinking around the given point, the submersion theorem provides
coordinates
\[
            (z,q_1,\ldots,q_r,u_1,\ldots,u_{d-r-1}).
\]
The coordinate $z$ is retained because $dq$ already has rank $r$ on its
zero divisor.  On the punctured neighborhood, the kernel inclusion and
the rank bound imply
$\rank(d\ell)=r$ and $\ker(dq)=\ker(d\ell)$.
Fix $q$ and $u$ and make one positive turn in the $z$ coordinate.
Along every lift of that loop, the derivative of $\ell$ vanishes, so its
value is constant.  Its endpoints are related by $T$, which consequently
fixes the initial line.  Varying the fixed coordinates proves the claim
on an open subset.  The projective equality $T\ell=\ell$ then extends by
the identity theorem to the connected universal cover.  For $r=0$, the
rank bound makes $\ell$ locally constant and the same loop argument
applies with no $q$ coordinates.
\end{proof}

\begin{lemma}[Boundary rank drop]\label{p1logalt:boundary-drop}
Let $Z$ be smooth projective, let $D$ be a reduced simple normal crossing
divisor, and let a polarizable integral pure variation on $Z\setminus D$
have a top Hodge line satisfying Condition~\eqref{p1logalt:faithful-line}.  Write $A_1$
for its rational extension and $r=\nu(A_1)$.  If $D_i$ is a component
whose transverse monodromy is not the identity, then
\[
 \nu(A_1|_{D_i})<r.
\]
\end{lemma}

\begin{proof}
If the boundary rank did not drop, suitable single-valued functions of
the moving line would have the same differential kernel as that line.
A transverse loop in one of their fibers would then force the monodromy
to fix the moving line on an open set, contradicting
Condition~\eqref{p1logalt:faithful-line} after analytic continuation.

Work at a general point of $D_i$, with coordinates $(z,s)$ on
$\Delta\times\Delta^{d-1}$ and $D_i=\{z=0\}$.  Write
$T=T_sT_u$ for the Jordan decomposition of the transverse monodromy and
put $N=\log T_u$.  Choose $e>0$ with $T_s^e=1$ and set $z=t^e$.
In flat coordinates the nilpotent orbit extension supplies a holomorphic
line map $\Psi$ on the root polydisk such that
\begin{equation}\label{p1logalt:untwist}
 \ell(t,s)=
 \exp\left(\frac{e\log t}{2\pi\sqrt{-1}}N\right)\Psi(t,s),
 \qquad
 \Psi(\zeta t,s)=T_s\Psi(t,s),
 \quad \zeta=\exp(2\pi\sqrt{-1}/e).
\end{equation}
These are equalities of lines; $\ell$ is evaluated on the logarithmic
cover when necessary.

Let $W=W(N)$, centered at the weight of the variation.  On the boundary
there is a unique index $k$ in which the top line has a nonzero
weight-graded image.  The map $N$ is a strict morphism of limiting mixed
Hodge structures of type $(-1,-1)$.  Since $F^{p+1}=0$, it gives
\begin{equation}\label{p1logalt:survives}
 F^p\cap NV=N(F^{p+1})=0.
\end{equation}
Thus the limiting line survives in $C=V/NV$.

We also need to preserve its derivative rank.  The primitive decomposition
of the fixed graded local system identifies
\[
 \operatorname{Gr}^W_k V
 =P_k\oplus N\operatorname{Gr}^W_{k+2}V,
 \qquad
 P_k\xrightarrow{\ \sim\ }\operatorname{Gr}^W_k C.
\]
Here $k$ is at least the original weight: lower graded pieces are
entirely in the image of $N$ and cannot contain the top line.  The second
summand has zero $F^p$, so the graded moving top line lies in the flat
primitive summand $P_k$.  Its projective derivative rank is therefore
unchanged by the displayed isomorphism.  Projecting the boundary line
in $C$ further to its relevant graded piece recovers precisely this pure
line period map.  By Lemmas~\ref{p1logalt:extensions} and~\ref{p1logalt:rank},
its rank is
\[
 r_i=\nu(A_1|_{D_i}).
\]
This equality may be checked on a global adapted cover: restriction of
the pulled-back line to a component dominating $D_i$ is its generically
finite pullback, which preserves numerical dimension.

Equation~\eqref{p1logalt:untwist} makes $\Psi(0,s)$ an eigenline of $T_s$.
Its eigenvalue $\lambda$ is constant on a sufficiently small connected
boundary neighborhood.  Let $E=C_\lambda$ and let $q_0:V\to E$ be the
quotient followed by the eigenspace projection.  The projective map
\[
 Q(t,s)=\mathbb P(q_0)(\Psi(t,s))
\]
is defined near the chosen boundary point and has boundary rank at
least $r_i$.  Since $N$ vanishes on $C$, it also equals
$\mathbb P(q_0)(\ell(t,s))$.  Moreover $T_s$ acts by a scalar on $E$,
so $Q(\zeta t,s)=Q(t,s)$.  Hence $Q$ descends holomorphically to the
original $(z,s)$-polydisk, including $z=0$.  This descent uses a
projective map; a choice of a fractional scalar frame is unnecessary.

Suppose $r_i\ge r$.  Choose $r$ affine coordinates of $Q$, denoted
$q_1,\ldots,q_r$, whose differentials are independent along $D_i$.
The map $q=(q_1,\ldots,q_r)$ factors locally through $\ell$ on the
punctured chart.  Its differential therefore satisfies the kernel
inclusion in Lemma~\ref{p4int:transverse-loop}.  The rank bound for
$\ell$ is $r$ by Lemma~\ref{p1logalt:rank}.  The transverse-loop lemma
shows that $T$ fixes the moving highest line everywhere on the connected
universal cover.  Condition~\eqref{p1logalt:faithful-line} gives
$T=\id$, contrary to the choice of the divisor.  If $r=0$, the same
lemma applies with no coordinate functions.  If $r>d-1$, the supposed
inequality $r_i\geq r$ is already impossible.
\end{proof}

\subsection{Removing the boundary from the adjoint class}

\begin{lemma}[A boundary subtraction estimate]\label{p1logalt:volume}
Let $Z$ be smooth projective of positive dimension $d$, let $A_1$ be a
nef rational divisor, and let $J=\sum_iD_i$ be reduced with smooth
components.  Suppose $K_Z+J$ is big and
\[
 \nu(A_1|_{D_i})<r:=\nu(A_1)\quad\text{for every }i.
\]
Then $K_Z+aA_1$ is big for every sufficiently large rational $a$.
\end{lemma}

\begin{proof}
If $r=0$, the restriction hypothesis forces $J=0$, so $K_Z$ is already
big.  Suppose $r>0$.  By Kodaira's lemma write
\[
 K_Z+J\sim_{\Q}H+E,\qquad H\text{ ample},\quad E\ge0,
\]
and set $B_a=H+aA_1$ for rational $a\ge0$.  Choose one ample integral
divisor $C_0$ such that both $C_0$ and $C_0+J$ are globally generated,
and put $H'=H+C_0$.  This choice is independent of $a$ and of all
section degrees below.

For every sufficiently divisible $m>0$, filter by $kJ$ for $0\le k<m$.
The restriction sequences give
\[
 h^0\bigl(Z,m(B_a-J)\bigr)
 \ge h^0(Z,mB_a)
 -\sum_{k=0}^{m-1}h^0\bigl(J,(mB_a-kJ)|_J\bigr).
\]
Since $J$ is reduced, its structure sheaf injects into
$\bigoplus_i\cO_{D_i}$.  Also
\[
 mC_0+kJ=(m-k)C_0+k(C_0+J)
\]
is globally generated.  On each $D_i$, multiplication by a section
which is not identically zero bounds the restricted summand by
$h^0(D_i,m(H'+aA_1)|_{D_i})$.  We obtain the exact estimate
\begin{equation}\label{p1logalt:uniform-sections}
 h^0\bigl(Z,m(B_a-J)\bigr)
 \ge h^0(Z,mB_a)
 -m\sum_i h^0\bigl(D_i,m(H'+aA_1)|_{D_i}\bigr).
\end{equation}
Fix $a$ and divide by $m^d/d!$.  Asymptotic Riemann--Roch for the ample
divisors involved, along divisible $m\to\infty$, yields
\begin{equation}\label{p1logalt:uniform-volume}
 \vol(B_a-J)
 \ge B_a^d-d\sum_i(H'+aA_1)^{d-1}\cdot D_i.
\end{equation}
No uniform Riemann--Roch remainder in $a$ is needed: the exact
inequality in Equation~\eqref{p1logalt:uniform-sections} holds for every $a$, and the
limit is taken with $a$ fixed.

The first polynomial on the right of
Equation~\eqref{p1logalt:uniform-volume} has degree $r$, with positive leading
coefficient $\binom dr A_1^rH^{d-r}$.  Every restriction polynomial has
degree at most $\nu(A_1|_{D_i})\le r-1$.  The right side is therefore
positive for every sufficiently large $a$.  Finally
$K_Z+aA_1\sim_{\Q}(B_a-J)+E$ is big.
\end{proof}

\begin{proof}[Proof of Theorem~\ref{p1logalt:positivity}]
Apply Lemmas~\ref{p1logalt:kernel} and~\ref{p1logalt:descent}.  We obtain a
variation on a dense open of a smooth projective $Z$ satisfying
Condition~\eqref{p1logalt:faithful-line}, with generically immersive full period map,
and rational Hodge extension $A_1$ satisfying
$b\pi^*L\sim_{\Q}h^*A_1$.  Resolve the boundary so that it is simple normal crossing.  Pullback
compatibility retains the line identity.  Delete every boundary component with identity transverse
monodromy, and call the remaining reduced divisor $J$.

The variation extends to $Z\setminus J$.  Indeed, at a point lying only
on deleted components all local monodromies are the identity, and the
nilpotent orbit extension with all monodromy logarithms zero is a pure
polarized extension.  The full period map is still generically immersive.
Thus \cite[Theorem~1.1 and Remark~1.2(1)]{BrunebarbeCadorel} applies
to the polarized variation on $Z\setminus J$ and gives the bigness of
$K_Z+J$; its hypothesis is immersivity at one point, and it imposes no
extra condition on boundary monodromy.
Lemma~\ref{p1logalt:boundary-drop} applies to every component of $J$, and
Lemma~\ref{p1logalt:volume} proves that $K_Z+aA_1$ is big for all sufficiently
large rational $a$.  When $\dim Z=0$ the assertion holds by convention.
Set $A=A_1/b$.  This is nef, satisfies
$\pi^*L\sim_{\Q}h^*A$, and has the same asserted adjoint positivity
after rescaling the large parameter.
\end{proof}

\begin{remark}\label{p1logalt:no-abundance}
The construction permits $\nu(A)<\dim Z$.  In particular, the algebraic
map $h$ is obtained from the full period map of an integral tensor
replacement; its fibers are not asserted to be the leaves of the line
period map.  Neither abundance nor semiampleness of $A$ occurs in the
argument.
\end{remark}

\section{Chern intersections and boundary limits of a highest Hodge line}
\label{p1varhd:section}

We prove the Chern-intersection formula deferred in
Lemma~\ref{p1logalt:rank}, and then give a different proof of the
boundary restriction estimate in Lemma~\ref{p1logalt:boundary-drop}.  Its two ingredients are a Chern-intersection
formula with arbitrary smooth closed test forms and a family of Thom
forms concentrating on a boundary component.  The limiting Hodge norm
computes the tangential curvature at a general boundary point; a product
Poincar\'e bound controls the crossings and permits passage to the
restriction intersection number.

Throughout this section, $R$ is a smooth connected projective complex
variety, $D_R$ is a reduced simple-normal-crossing divisor, and
$\mathbb V$ is a polarizable integral pure variation on $R\setminus D_R$.
A highest Hodge line is its highest nonzero filtration piece $F^p$, with
$F^{p+1}=0$ and $\rank F^p=1$.  We denote its rational parabolic
extension by $A$.  The parabolic convention includes finite-monodromy
weights.  Its line period map is the map from the universal cover of
$R\setminus D_R$ to the projective space of a fixed flat reference fiber
which remembers $F^p$.  Its rank need not equal the full-period rank.
For $n=\dim R$ we use
\[
 \nu(A)=\max\{j\in\{0,\ldots,n\}:
                   c_1(A)^j\cdot H^{n-j}>0\},
\]
where $H$ is any ample class and $A$ is nef.  In particular, this
numerical dimension is zero on a point.  Normalize
$\ddc=\frac{i}{2\pi}\partial\bar\partial$.

Lemma~\ref{p1logalt:extensions} supplies nefness and compatibility of
the rational extension with positive tensor powers and log-pair
pullbacks whose interiors map into the original interior. The
Chern-intersection formula below uses these extension properties,
but neither full-period descent nor a boundary rank-drop result.
For strict rank loss we also use the integral tensor replacement
of Lemma~\ref{p1logalt:kernel} and the transverse-loop criterion of
Lemma~\ref{p4int:transverse-loop}. The replacement is faithful on
the moving highest line: a monodromy element fixing that line at
every point of the connected universal cover is the identity.
The Thom-form calculation computes restriction intersections from
tangential limiting curvature. It uses the local graded Hodge data,
without identifying the boundary restriction with a globally extended
graded line as in Section~\ref{p1logalt:section}.

\subsection{Curvature and intersection powers}

\begin{lemma}[Chern intersections with closed test forms]\label{p1varhd:curvature}
Let $A$ be the rational parabolic highest line of an integral
polarized pure variation on the complement of a reduced
simple-normal-crossing divisor in a smooth projective variety $R$.
Then
\begin{equation}\label{p1varhd:numrank}
 \nu(A)=\text{generic differential rank of its projective line map}
\end{equation}
provided $A$ is nef.  In the unipotent case, if $\alpha$ is the Chern
form of its Hodge metric on $R\setminus D_R$, then for every
$0\leq j\leq n=\dim R$ and every smooth closed form $\gamma$ of type
$(n-j,n-j)$, the integral
$\int_{R\setminus D_R}\alpha^j\wedge\gamma$ is absolutely convergent
and equals $c_1(A)^j\cdot[\gamma]$.
\end{lemma}

\begin{proof}
We may compute after a generically finite smooth alteration on
which all boundary monodromies are unipotent.  To obtain one, pass
to a finite level cover with neat monodromy, compactify, and resolve.
Quasi-unipotence then gives unipotence, including for the commuting
local monodromies on the resolved boundary.  The parabolic pullback
property of Lemma~\ref{p1logalt:extensions} identifies the new line
with the pullback of $A$.
Generically finite pullback preserves nef numerical dimension by
the intersection and projection formulas.  It also preserves the
generic differential rank in characteristic zero.  We therefore
work on the unipotent model, retaining the notation $R,A$.

Let $\alpha$ be the semipositive Chern form of the highest-line
Hodge metric on the open part. Griffiths' curvature formula
\cite[Theorem~5.2, Formula~5.3, and Formula~6.18]{GriffithsPeriodsIII} says
that $\alpha(v,\bar v)$ is a positive fixed multiple of the squared
norm of the projective highest-line differential in direction $v$.
In particular, its rank is the rank in \eqref{p1varhd:numrank}.

Consider a crossing chart with coordinates $z_1,\ldots,z_n$ and
boundary $z_1\cdots z_c=0$.  Put
$\rho_i=-\log|z_i|^2$ and shrink the chart so that $\rho_i>e$.
The horizontal Schwarz estimate, applied separately to the
coordinate disks, gives
\begin{equation}\label{p1varhd:poincare}
 0\leq\alpha\leq C\Omega,
 \qquad
 \Omega=
 \sum_{i=1}^{c}\frac{i\,dz_i\wedge d\bar z_i}
                         {|z_i|^2\rho_i^2}
 +\sum_{i=c+1}^{n}i\,dz_i\wedge d\bar z_i.
\end{equation}
Here and below harmless positive normalization constants are
absorbed into $C$.  The coordinate estimate follows from the
curvature of the horizontal period metric and the singular
Ahlfors--Schwarz lemma; see \cite[\S 2.1, Proposition~2.4 and
Lemma~2.5]{BrunebarbeCadorel}.  Positivity and the Hermitian-form Cauchy--Schwarz
inequality control the off-diagonal entries and give the uniform
product estimate \eqref{p1varhd:poincare}, even where the period
differential is degenerate.

By the nilpotent-orbit theorem
\cite[Theorem~4.12]{Schmid}, a nonvanishing top extension frame,
written in flat coordinates on a bounded angular sector, is
\begin{equation}\label{p1varhd:frame}
 v(z)=\exp\left(\sum_{i=1}^{c}
               \frac{\log z_i}{2\pi i}N_i\right)u(z),
\end{equation}
where the $N_i$ commute and $u$ is holomorphic and nonvanishing.
If $\beta$ is the curvature of a fixed smooth metric on $A$, write
$\alpha-\beta=\ddc\varphi$.  The ratio of the two metrics satisfies
\begin{equation}\label{p1varhd:loglog}
 |\varphi(z)|\leq C\left(1+\sum_{i=1}^{c}\log\rho_i\right).
\end{equation}
This is the logarithmic growth estimate for canonical Hodge frames;
see \cite[Theorem~5.1]{CattaniKaplanDegenerating}.  Here the Schwarz
estimate also gives a direct proof of the bound needed for integration.  The differential of the logarithmic Hodge norm of any
flat vector is bounded by a constant times the invariant period
metric, with a constant independent of the point by homogeneity.
Travel from a fixed point along coordinate paths in a bounded
angular sector.  Their product-Poincar\'e length is bounded by
$C(1+\sum_i\log\rho_i)$, so the Hodge metric and its inverse in a
flat frame grow at most as products of powers of the $\rho_i$.
The exponential in \eqref{p1varhd:frame} and its inverse are polynomials
in the logarithms, since the $N_i$ are commuting nilpotent
operators.  Finally, $u$ has ordinary norm bounded above and below
on a smaller chart.  These three observations give the upper and
lower logarithmic bounds in \eqref{p1varhd:loglog}.  We have therefore
obtained the two controls needed for the cutoff argument:
\begin{equation}\label{p1logalt:metric-bounds}
 0\leq\alpha\leq C\Omega,
 \qquad |\varphi|\leq C\Bigl(1+\sum_i\log\rho_i\Bigr).
\end{equation}

Chern-current representation for canonical Hodge extensions is
established in \cite[Theorem~5.2]{CattaniKaplanDegenerating}. We now
justify the highest-line intersection formula with closed test forms,
including the absence of boundary contributions.  For $0\leq j\leq n$ and every
smooth closed form
$\gamma$ of type $(n-j,n-j)$,
\begin{equation}\label{p1varhd:chern}
 \int_{R\setminus D_R}\alpha^j\wedge\gamma
       =c_1(A)^j\cdot[\gamma].
\end{equation}
The case $j=0$ is immediate, so assume $j>0$.  Choose a smooth
function $\chi$ that is one on $(-\infty,1]$ and
zero on $[2,\infty)$, and use products of
$\chi((\log\rho_i)/M)$ to cut off the boundary.  The $\rho_i$ can
be globalized using smooth metrics on the boundary divisor
bundles.  Denote the resulting cutoff by $\chi_M$.  Its second
derivatives are uniformly bounded in the product-Poincar\'e metric,
and their support escapes to the boundary as $M\to\infty$.

Since all curvature forms are closed, on the open part we have
\[
 \alpha^j-\beta^j
 =\ddc\left(\varphi
       \sum_{k=0}^{j-1}\alpha^k\wedge\beta^{j-1-k}\right).
\]
Multiplying by $\chi_M\gamma$ and integrating by parts makes the
difference of the integrals equal to
\[
 \int\varphi\,\ddc\chi_M\wedge
       \sum_{k=0}^{j-1}\alpha^k\wedge\beta^{j-1-k}\wedge\gamma.
\]
Its absolute value is bounded on escaping tails by a constant times
\[
 \left(1+\sum_i\log\rho_i\right)\Omega^n.
\]
This is integrable: in a cusp variable the relevant radial
integrals are bounded by
$\int^{\infty}(1+\log\rho)\rho^{-2}\,d\rho$.
Thus the error tends to zero.  The same domination permits removal
of $\chi_M$ from the two curvature integrals, proving
\eqref{p1varhd:chern}.  Rational bundles are handled by first taking an
integral tensor power and then dividing the forms and identities
by that power.

Let $r$ be the generic rank of the projective line map.  Then
$\alpha^{r+1}=0$ on the open part, whereas $\alpha^r$ is nonzero
and positive on a nonempty open subset.  Apply \eqref{p1varhd:chern}
with complementary powers of an ample K\"ahler form.  It gives
vanishing of the $(r+1)$st intersection power and strict positivity
of the $r$th one.  The nef intersection characterization proves
\eqref{p1varhd:numrank}.
\end{proof}

\subsection{Rank loss at a monodromy divisor}

\begin{lemma}[Boundary rank drop by a Thom-form calculation]\label{p1varhd:boundary}
In the setting of Lemma~\ref{p1varhd:curvature}, suppose that monodromy
is faithful on the moving highest line.  If $D_i$ is a boundary
component with nonidentity transverse monodromy, then
\begin{equation}\label{p1varhd:strictdrop}
             \nu(A|_{D_i})<\nu(A).
\end{equation}
\end{lemma}

\begin{proof}
Put $r=\nu(A)$.  We first compute a bound for the numerical
dimension on a boundary component of a unipotent model.  This
computation is unchanged on a component dominating $D_i$ under a
generically finite alteration, since restriction of the pulled-back
line is the generically finite pullback of $A|_{D_i}$.

\paragraph{The limiting highest line.}
Near a general point of a unipotent boundary component write the
coordinates as $(t,s)$, with divisor $t=0$ and monodromy logarithm
$N$.  Let $w$ be the weight and $p$ the largest filtration index.
The limiting mixed Hodge structure is polarized by $N$
\cite[Theorem~6.16]{Schmid}; its graded structures vary as polarized
variations along the smooth boundary stratum
\cite[Proposition~2.10 and the following discussion]{CattaniKaplanDegenerating}.
The limiting $F^p$ has dimension one and $F^{p+1}=0$.  Consequently
it consists of one Deligne summand $I^{p,q}$.  Write
$p+q=w+l$.  In the primitive decomposition for $N$, a nonzero
$N^k$-image term with $k>0$ would have to come from $F^{p+k}$,
which is zero.  The line is therefore primitive, $l\geq0$, and,
for a nonzero vector $u_0(s)=u(0,s)$ spanning it,
\[
       N^lu_0(s)\ne0,\qquad N^{l+1}u_0(s)=0.
\]
Shrink the boundary stratum so that its Hodge type, and hence $l$,
is fixed.

For a highest Hodge vector in the pure variation, the Hodge norm is
the flat polarization pairing with its conjugate, multiplied by the
fixed factor determined by its Hodge type.
Substituting the exponential frame into that pairing gives a
polynomial in $y=-\log|t|$.  Primitive polarization identifies its
leading term as a positive constant times $y^l$ times the squared
norm $h_l(s)$ of the primitive graded line.  Replacing $u(t,s)$
by $u_0(s)$ introduces errors $O(|t|)$ times powers of $y$;
the same holds after tangential differentiation, by holomorphicity
of $u$.  It follows, locally uniformly with tangential derivatives,
that
\begin{equation}\label{p1varhd:limitnorm}
 \log\|v(t,s)\|^2
      =l\log y+\log h_l(s)+c+O(y^{-1})+O(|t|y^M)
\end{equation}
for some $M$ and a constant $c$.  Thus the tangential curvature
converges to the curvature of the primitive graded highest line.
If $q_i$ denotes the generic rank of its projective differential,
the limiting tangential curvature has rank at most $q_i$.

\paragraph{The restriction intersection and crossing points.}
We claim that
\begin{equation}\label{p1varhd:restrictionrank}
                    \nu(A|_{D_i})\leq q_i.
\end{equation}
The following estimate proves the claim without assuming uniform
convergence of \eqref{p1varhd:limitnorm} at crossings.  Fix an ample
K\"ahler form $\omega$ and an integer $j$ with
$q_i<j\leq n-1$.  Choose a smooth metric on $\cO(D_i)$, write its
curvature as $\beta_i$, and put
$\tau=\log\|s_{D_i}\|^2$.  Let $g$ be a fixed smooth convex
function with $g'=0$ on $(-\infty,0]$ and $g'=1$ on
$[\log4,\infty)$, and set
$f_\epsilon(x)=g(x-\log\epsilon^2)$.  Then
\begin{equation}\label{p1varhd:thom}
 \eta_\epsilon
  =\beta_i+\ddc f_\epsilon(\tau)
  =(1-f_\epsilon'(\tau))\beta_i
     +\frac{i}{2\pi}f_\epsilon''(\tau)
                   \partial\tau\wedge\bar\partial\tau
\end{equation}
is a smooth closed representative of $c_1(\cO(D_i))$.
Here the normalization is $\ddc\log|t|^2=[t=0]$.
The function $f_\epsilon(\tau)$ is constant near $D_i$, so it
extends smoothly there.  Outside $\|s_{D_i}\|\leq2\epsilon$,
the two curvature terms cancel and $\eta_\epsilon=0$.

On a crossing chart put
\[
 \lambda=\frac{i\,dt\wedge d\bar t}{|t|^2},\qquad
 \rho=-\log|t|^2,\qquad
 \Omega=\frac{\lambda}{\rho^2}+\Omega_s,
 \qquad T=\alpha^j\wedge\omega^{n-1-j}.
\]
The tangential form $\Omega_s$ is a product cusp metric, with
Euclidean summands in the remaining coordinates.  It has finite
volume.  The bound \eqref{p1varhd:poincare} gives
$0\leq T\leq C\Omega^{n-1}$.
Write $\tau=\log|t|^2+\psi$ with $\psi$ smooth, and put
$L_\epsilon=|\log\epsilon|$.  The smooth first term of
\eqref{p1varhd:thom} contributes at most
\[
 C\int_{|t|\leq C\epsilon}\Omega^n
                     =O(L_\epsilon^{-1}).
\]
For the second term write $\partial\tau=a_0+b_0$, with
$a_0=dt/t$ and $b_0=\partial\psi$.  Its support lies in
$c\epsilon\leq|t|\leq C\epsilon$, with constants independent
of $s$.  On this shell, positivity gives
\begin{align*}
 \int i a_0\wedge\bar a_0\wedge T&=O(1),\\
 \int i b_0\wedge\bar b_0\wedge T
       &\leq C\int_{c\epsilon\leq|t|\leq C\epsilon}\Omega^n
         =O(L_\epsilon^{-2}).
\end{align*}
The first estimate follows because wedging with $\lambda$
selects only the tangential part of $T$.  For the second, a smooth
positive form dominating $i b_0\wedge\bar b_0$ is bounded by
$C\Omega$.  Cauchy--Schwarz for the positive form $T$ now bounds
the mixed terms by $O(L_\epsilon^{-1})$.
The bounded coefficient $f_\epsilon''$ does not change these
estimates.

It remains to consider the pure normal-shell term.  Set
$t=\epsilon u$.  Its cutoff coefficient becomes
\[
 g''\bigl(\log|u|^2+\psi(\epsilon u,s)\bigr),
\]
whose support is contained in a fixed annulus, since $\psi$ is
bounded.  The absolute integrand is bounded, independently of
$\epsilon$, by
\[
 C\,\frac{i\,du\wedge d\bar u}{|u|^2}
            \wedge\Omega_s^{n-1}.
\]
This is integrable on that annulus times the tangential chart,
including all tangential divisor crossings.  For almost every $s$,
\eqref{p1varhd:limitnorm} gives a tangential curvature limit of rank
less than $j$, so the tangential $j$th power tends to zero.
Dominated convergence makes the pure normal-shell integral tend
to zero.  The estimates are uniform under truncating the
tangential cusp variables, so no mass is lost when those
truncations are removed.  Choose a finite partition of unity subordinate to crossing charts.
Multiplication by its bounded smooth functions preserves each absolute
estimate.  Summing the local estimates proves the same conclusion on
all of $D_i$; no integration by parts is performed after inserting this
partition.

For each fixed $\epsilon$, apply \eqref{p1varhd:chern} with
$\gamma=\omega^{n-1-j}\wedge\eta_\epsilon$.  Its cohomology
class is independent of $\epsilon$.  The preceding limit gives
\[
 c_1(A|_{D_i})^j\cdot[\omega|_{D_i}]^{n-1-j}=0.
\]
Nefness proves \eqref{p1varhd:restrictionrank}.  If $q_i=n-1$, that
inequality is automatic.  This also explains why the generic
limit calculation alone would not be sufficient: the uniform
product bound is what controls the crossing locus.

\paragraph{A fixed quotient and transverse loops.}
Return to the original boundary component, with transverse
monodromy $T_s\exp N$, where $T_s$ has finite order.  Work first
in a transverse root coordinate killing $T_s$.  Set
$C_N=\operatorname{coker}N$, endowed with the quotient of the
fixed monodromy weight filtration.  The primitive limiting line
of weight $k=w+l$ projects nontrivially into $W_kC_N$.
The fixed linear quotient
\begin{equation}\label{p1varhd:fixedquotient}
              W_kC_N\longrightarrow\operatorname{Gr}^W_k C_N
\end{equation}
recovers its primitive graded line.  Indeed, in the Lefschetz
decomposition of $\operatorname{Gr}^W_k$, quotienting by the
image of $N$ removes exactly the summands that are positive powers
of $N$ applied to higher primitives.  The induced quotient is the
primitive part of weight $k$.  In particular, projective
differential rank cannot increase under \eqref{p1varhd:fixedquotient}.
No choice of a moving primitive splitting is involved.

The limiting line in $C_N$ lies in one $T_s$-eigenspace.  Compose
the quotient by $N$ with the fixed projection onto this eigenspace.
Near the chosen general boundary point this gives a nonzero
holomorphic projective map from the untwisted highest line.
Under deck rotation the untwisted line is acted on by $T_s$;
on the projected eigenspace that action is scalar.  The projective
map therefore descends to the original chart.  On the punctured
chart it factors through the moving highest-line map, since the
exponential of $N$ disappears in $C_N$.  Its boundary differential
rank is at least the rank $q_i$ of the primitive graded line.

Suppose $\nu(A|_{D_i})\geq r$.  By
\eqref{p1varhd:restrictionrank}, the projected map has at least $r$
independent tangential differentials at a general boundary point.
Choose $r$ affine components $g_1,\ldots,g_r$ with that property.
Their differentials are independent along the boundary at that point,
and the map $g=(g_1,\ldots,g_r)$ extends holomorphically over the original
chart.  On the punctured chart it factors through the highest-line map.
Lemma~\ref{p1varhd:curvature} gives the latter rank at most $r$.
The hypotheses of the transverse-loop criterion
(Lemma~\ref{p4int:transverse-loop}) are therefore satisfied in the
original transverse coordinate, before the finite root cover.  It
follows that the original transverse monodromy fixes the moving line
on the connected universal cover.  This contradicts faithfulness and
its nonidentity.  When $r=0$, use the same lemma with an empty list of
functions; when $r>n-1$, the supposed inequality
$\nu(A|_{D_i})\geq r$ is impossible on dimensional grounds.  This proves
Equation~\eqref{p1varhd:strictdrop}.
\end{proof}


\subsection{Connection with the integral projective argument}

The hypotheses required for the preceding calculation arise from the
integral tensor replacement and full-period descent in
Lemmas~\ref{p1logalt:kernel} and~\ref{p1logalt:descent}.
On the resulting projective base, Lemma~\ref{p1varhd:boundary} gives the
same strict inequality as Lemma~\ref{p1logalt:boundary-drop}, with
the boundary intersection justified directly by Thom forms.  The
boundary subtraction estimate of Lemma~\ref{p1logalt:volume} then
applies to the union $J$ of the divisors with nonidentity transverse
monodromy.  The full-period map supplies the bigness of $K_R+J$ as in
the proof of Theorem~\ref{p1logalt:positivity}; the highest-line
rank supplies the strict decrease in the degree of the restriction
polynomials.  Thus the two different period ranks retain their separate
roles in adjoint positivity.

\begin{remark}[Tensor replacement on a restricted locus]
\label{p4int:restricted-tensor}
The same local results also give the curve-constancy part
of the highest-line properties used in such comparisons.  On a smooth
complete curve, degree zero gives zero integral of the semipositive
Hodge curvature by Equation~\eqref{p1varhd:chern}, after an adapted
cover if necessary.  The curvature is therefore zero, and the highest
line is projectively constant in a flat reference.  If the highest line
is projectively constant on a connected smooth algebraic locus, apply
Lemma~\ref{p1logalt:kernel} to the restricted variation itself.
Use its own monodromy image and isotypic decomposition: restricting
an ambient replacement need not preserve faithfulness, because the
restricted line image may have a larger projective stabilizer.
Every monodromy element of the replacement fixes its constant highest
line, so faithfulness makes that monodromy trivial.  Since the new line
is a positive tensor power of the original line, the original
highest-line character on that locus has finite image.  The tensor
exponent here is allowed to depend on the restricted variation.
\end{remark}

\begin{remark}\label{p1varhd:further-models}
The adjoint conclusion is preserved under further smooth projective
birational modification $\pi:R'\to R$, together with a compatible
modification of the source.  The parabolic pullback property preserves
the line-bundle identity and nefness, while
\[
 K_{R'}+b\pi^*A
   =\pi^*(K_R+bA)+E_\pi,\qquad E_\pi\geq0,
\]
preserves bigness.  The calculations here concern the whole rank-one
highest step of a polarizable integral pure variation on a projective
pair; they assert no equality between highest-line rank, full-period
rank, and whole-fiber birational variation.
\end{remark}

\bibliographystyle{amsplain}
\bibliography{references}
\end{document}
